\documentclass{article}

\PassOptionsToPackage{numbers,sort&compress}{natbib}
\usepackage[final,eandd]{neurips_2026}

\usepackage[utf8]{inputenc}
\usepackage[T1]{fontenc}
\usepackage[hypertexnames=false,hidelinks]{hyperref}
\usepackage{url}
\usepackage{booktabs}
\usepackage{amsmath}
\usepackage{amsfonts}
\usepackage{amssymb}
\usepackage{bbm}
\usepackage{nicefrac}
\usepackage{xcolor}
\usepackage{graphicx}
\usepackage{array}
\usepackage{algorithm}
\usepackage{algpseudocode}
\usepackage{cleveref}
\usepackage{multirow}

% ---- placeholder macros -----------------------------------------
\newcommand{\TODO}[1]{{\color{red}\textbf{[TODO: #1]}}}
\newcommand{\samethanks}{\footnotemark[\value{footnote}]}
\newcommand{\todofig}[1]{\fbox{\parbox[c][0.6in][c]{0.95\linewidth}{\centering
  \small\color{red} placeholder figure: \texttt{#1}}}}

\title{MatCurvs: Article Real-Coordinate Curve Extraction\\
for Agent-Ready Materials Reasoning}

\author{%
  Liang Yin\\
  School of Advanced Interdisciplinary Sciences,\\
  University of Chinese Academy of Sciences\\
  \texttt{yinliang22@mails.ucas.ac.cn}
  \And
  Songlin Yu\thanks{Corresponding author.}\\
  University of Chinese Academy of Sciences\\
  \texttt{yusonglin22@mails.ucas.ac.cn}
  \And
  Jianjun Liu\samethanks\\
  Shanghai Institute of Ceramics,\\
  Chinese Academy of Sciences\\
  \texttt{jliu@mail.sic.ac.cn}
}

\begin{document}
\maketitle
\thispagestyle{plain}  % override .sty's empty: show page-1 page number

% --- Teaser figure: full pipeline (PDF -> real-coord curves).
% Placement: AFTER title+authors, BEFORE abstract on page 1.
% Using fig_doc_pipeline.png as the teaser asset (existing 3-stage
% schematic). Replace with fig_full_pipeline.png if a dedicated
% PDF-to-real-coord overview figure is delivered.
\begin{figure}[!h]
\centering
\includegraphics[width=\linewidth]{figures/fig_doc_pipeline.png}
\caption{\textbf{The MatCurvs end-to-end pipeline.} From a raw
materials-science PDF (left) through subfigure detection and a
$3{\times}10$ characterisation taxonomy (L0), pixel-polyline GT for
$\mathrm{LR}_{\theta}$ (L1), Mask2Former + 5D LAB-$xy$ $k$-means
extraction with VLM-judge curation (L2), to real-coordinate
$(x_{\mathrm{real}}, y_{\mathrm{real}})$ curves feeding agentic
numerical-readout and per-family domain-feature reasoning.}
\label{fig:teaser}
\end{figure}

% ================================================================
\begin{abstract}
Materials-science articles communicate experimental results almost
entirely through line plots, but the agents that read this
literature today consume those plots only as VLM-described prose.
We measure two failures. \textbf{(1) Pixel anchoring}: four hosted
VLMs (GPT-5.4, Claude Opus 4.6, Gemini-3.1-pro-preview, qwen-vl-max)
cap at $\mathrm{LR}_{0.5}{\le}1.2\%$, while a local
Mask2Former${+}$5D LAB-$xy$ $k$-means baseline reaches $45.8\%$ ---
a $\mathbf{38\times}$ gap. \textbf{(2) Numerical reasoning}: on
$14{,}740$ deterministic questions over $3{,}817$ panels, image-
only VLMs miss the answer by a median $15\text{--}54\%$; injecting
the extracted real-coordinate $(x,y)$ data into the prompt drops
error to $4.3\%$ for the same VLM, and an image-blind text LLM
with data alone reaches $4.2\%$. \emph{The image becomes redundant
once the model has the numbers.} Direct readouts ($y$ at $x_0$,
EIS $Z'$ intercept) drop from $99.8\text{--}100\%$ to
$0\text{--}1.5\%$ ($\mathbf{67\times}$ to $\mathbf{\infty\times}$);
the stacked-spectra intensity-offset convention drops from $99\%$
to $3\%$ ($\mathbf{33\times}$). We release \textbf{MatCurvs}
(grounded in $55{,}763$ articles: $1{,}375{,}165$ classified
subfigures, $204$ manually pixel-polylined panels, $8{,}026$
verified real-coordinate curves over $3{,}817$ panels, $14{,}740$
reasoning questions), the \textbf{MatDeplot} local extraction
pipeline, and a four-method evaluation pinpointing which
agentic-reasoning tasks pixel-faithful extraction unlocks.
\end{abstract}

% ================================================================
\section{Introduction}
\label{sec:intro}

\paragraph{Pain point 1: figures in materials-science papers are composite.}
A typical materials-science paper rarely shows one experiment per
figure. The same figure routinely bundles an XRD pattern, a
Raman spectrum, an SEM micrograph, a GCD curve, a CV scan, and a
process schematic into a single composite layout with no
machine-readable indication of which panel measures \emph{what} or
which axes carry which units (\Cref{fig:teaser}). An agent that
wants to reason over the article must first decompose this layout
into individual panels and classify each panel before any
numerical extraction is possible.

\paragraph{Pain point 2: the category space is large and quantitative content concentrates in line plots.}
Across the $55{,}763$ materials-science articles we processed, our
classifier resolves the panel space along a $6$-chart-type
$\times$ $3$-level $\times$ $10$-second-level characterisation
taxonomy: microscopy (SEM/TEM), schematics, computational images
(DFT/MD/FEM), and seven distinct line-plot families
(XRD, Raman, XPS, XAFS, EIS, GCD, CV) each with its own axis
convention and physical interpretation. The total subfigure count
is $\mathbf{1{,}375{,}165}$, of which $\mathbf{47.81\%}$
($\mathbf{657{,}428}$) are line plots. The line plots
disproportionately carry the paper's first-class numerical claims
(XRD peak positions $\to$ phase identification, Raman peak
positions and ID/IG ratio $\to$ disorder/graphitisation, GCD
plateau capacity $\to$ battery performance, XAFS edge position
$\to$ oxidation state). Get the line plots right and most of the
paper becomes machine-readable; get them wrong and downstream
reasoning is built on noise.

\paragraph{Pain point 3: today's only access path is a hosted VLM, and it cannot read the numbers.}
Today the only viable agentic access to those line plots is via
hosted vision--language models~\citep{gpt4,claude,gemini,qwenvl}.
VLMs describe a curve's overall shape with high prose-level
fluency (Inst-F1 ${\approx}0.93$) but they cannot anchor their
output to source pixels and they cannot recover axis-calibrated
$(x,y)$ values: at the dataset's $5$-px GT rasterisation, at most
$\mathbf{1.2\%}$ of VLM-predicted curves reach IoU${\ge}0.5$
against ground truth, uniformly across vendors. Prior local
foreground-segmentation pipelines~\citep{deeplabv3plus,unetplusplus,
resnest} silently collapse on dense-scatter spectra and emit zero
curves. \textbf{There is no widely-available conversion step that
turns a published chart image into the real-coordinate $(x,y)$
data an agent can compute on.} As a result, agents reading
materials literature today perform shape-level prose recognition,
not numerical reasoning --- and on the $14{,}740$-question agentic
benchmark we introduce in \S\ref{sec:reasoning}, image-only VLMs
miss the answer by a median of $15\text{--}54\%$ relative error.

\paragraph{Our work.}
We close the article-to-real-coordinate gap end-to-end. We release
\textbf{MatCurvs}, a three-tier benchmark grounded in the
$55{,}763$-article corpus: \textbf{(L0)} $\mathbf{1{,}375{,}165}$
classified subfigures with chart-type and two-level taxonomy;
\textbf{(L1)} $n{=}204$ subfigures ($n{=}90$ leakage-clean) with
manual per-curve pixel polylines; \textbf{(L2)} $8{,}026$
semi-automatic, VLM-judge-verified real-coordinate curves over
$3{,}817$ panels and $2{,}344$ source articles spanning Raman,
XAFS, XRD, GCD, CV, EIS and XPS. We release \textbf{MatDeplot}, a
modular local baseline (Mask2Former-q12~\citep{mask2former,mitb3}
with multi-scale${+}$flip TTA, a 5D LAB-$xy$ $k$-means separator
with foreground-fraction-triggered fallback, per-curve abstain) that
achieves $\mathrm{LR}_{0.5}{=}45.8\%$ ($\boldsymbol{38\times}$ over
the best VLM) while matching it on Inst-F1 ($0.942$ vs.\ $0.930$).
We release \textbf{MatCurvs-Reasoning}, a downstream evaluation
showing that real-coordinate curves enable agentic numerical readout
and per-family domain-feature extraction at instrument resolution
--- reasoning capabilities VLM-only pipelines cannot perform.

\paragraph{Contributions.}
\begin{itemize}\itemsep -2pt
\item An \textbf{article-to-real-coordinate pipeline} scaling from
$55{,}763$ raw articles to $\approx 8{\times}10^{3}$ verified
real-coordinate curves across seven characterisation families.
\item The \textbf{MatCurvs three-tier benchmark}
($1{,}375{,}165$\,/\,$204$\,/\,$8{,}026$ entries) curated under a
documented \emph{semi-automatic with VLM-judge verification}
annotation protocol with an inter-judge agreement audit
(\Cref{ssec:l2}).
\item A \textbf{baseline that closes the gap}: MatDeplot reaches
$\mathrm{LR}_{0.5}{=}45.8\%$ ($38\times$ over the best VLM), and
\textbf{MatCurvs-Reasoning} shows that the resulting real-coordinate
curves enable numerical point readout and per-family domain feature
extraction at instrument resolution --- tasks VLM-only pipelines fail.
\end{itemize}

\Cref{sec:related} situates this work;
\Cref{sec:formal} formalises real-coordinate chart reading and
$\mathrm{LR}_{\theta}$;
\Cref{sec:bench} describes MatCurvs and the annotation protocol;
\Cref{sec:method} details MatDeplot;
\Cref{sec:results} reports the pixel-anchor headline;
\Cref{sec:reasoning} reports MatCurvs-Reasoning;
\Cref{sec:limits} catalogues limitations.

% ================================================================
\section{Related Work}
\label{sec:related}

\paragraph{Chart derendering and chart-VQA benchmarks.}
Chart-to-table models~\citep{deplot,matcha,unichart,onechart,
chartassistant} target tabular charts (bar, pie, short line lists)
and produce zero recoverable polylines on dense XRD/Raman spectra
(\Cref{tab:method-vs-baselines}). Existing chart-VQA benchmarks
\citep{chartqa,plotqa,charxiv,chartbench} score
\emph{question-answering about} charts rather than \emph{numerical
extraction from} them; MatCurvs is, to our knowledge, the first
chart benchmark with per-curve pixel \emph{and} real-coordinate
annotations over materials-science chart families.

\paragraph{Hosted VLMs and dense segmentation.}
Frontier VLMs~\citep{gpt4,claude,gemini,qwenvl} match a curve's
shape in prose but cannot anchor to source pixels
(\S\ref{sec:results}). MatDeplot's foreground head is a
Mask2Former~\citep{mask2former} on a MiT-B3 encoder~\citep{mitb3}
with $q{=}12$ object queries, replacing the
DeepLabV3+/UNet++~\citep{deeplabv3plus,unetplusplus} on
ResNeSt~\citep{resnest} backbones used by prior local chart-FG
baselines, which we find dominant-mask-biased on dense scatter
(\Cref{tab:method-vs-baselines}).

\paragraph{Scientific literature mining and agent loops.}
WebPlotDigitizer~\citep{webplotdigitizer} digitises curves by hand;
ChemDataExtractor~\citep{chemdataextractor} and the text-mining
pipelines of \citet{kononova2019,olivetti2020} ingest text and
tables but not figures. ReAct/Toolformer-style agents
\citep{react,toolformer} amplify tool latency and cost --- our
abstention-supporting JSON output is designed for that regime.

% ================================================================
\section{Problem Formalization}
\label{sec:formal}

\paragraph{Inputs and outputs.}
A line-plot subfigure $I \in \mathbb{R}^{H \times W \times 3}$
depicts $K$ scientific curves within an axis-bounded data region.
A reader is asked, for each curve $i \in \{1,\ldots,K\}$, to produce
(a) a pixel set $\hat{C}_{i} \subset \{1,\ldots,H\}\times\{1,\ldots,W\}$
tracing the curve in image coordinates, and
(b) an ordered real-coordinate sequence
$\hat{S}_{i}=\bigl((x^{i}_{1},y^{i}_{1}),\ldots,(x^{i}_{n_i},y^{i}_{n_i})\bigr)$
in the chart's physical units, obtained via axis calibration.
Ground truth from L1 is the manual per-curve pixel polyline
$C^{*}_{j}$, rasterised at thickness $5$\,px (consistent with
published GT-line conventions). Ground truth from L2 is the
VLM-judge-verified real-coordinate sequence $S^{*}_{j}$.

\paragraph{Pixel anchoring: line recall $\mathrm{LR}_\theta$.}
Given prediction $\{\hat{C}_{i}\}$ and L1 ground truth
$\{C^{*}_{j}\}$, we compute pairwise IoU, solve the Hungarian
assignment $\sigma(\cdot)$, and define
\begin{equation}
  \mathrm{LR}_{\theta} \;=\;
  \frac{1}{K}\sum_{j=1}^{K}
  \mathbbm{1}\!\left[\mathrm{IoU}(C^{*}_{j},\hat{C}_{\sigma(j)})\ge\theta\right],
  \label{eq:lr}
\end{equation}
the fraction of GT curves whose matched prediction exceeds threshold
$\theta$ at the dataset's $5$-px rasterisation. We report $\theta\in
\{0.5, 0.8\}$ as the principal headline metric.

\paragraph{Permissiveness hierarchy.}
Three nested levels of agreement let a reader claim it has ``read''
a chart: \textbf{Shape} --- curves match morphology (Inst-F1, DTW)
without anchoring to pixels; \textbf{Instance} --- curves are
correctly partitioned into per-curve identities (right $K$, right
colour-to-curve assignment); \textbf{Pixel} --- curves overlap
source pixels at IoU${\ge}\theta$.
$\text{Pixel}\!\subset\!\text{Instance}\!\subset\!\text{Shape}$ as
classifier strength: any pixel-anchored prediction is also instance-
and shape-correct, but not vice versa. Hosted VLMs occupy
Shape-only; real-coordinate reading requires Pixel \emph{plus} axis
calibration.

\paragraph{Real-coordinate agreement.}
Beyond pixel anchoring, we measure real-coordinate agreement
between $\hat{S}_{i}$ and $S^{*}_{i}$ via dynamic-time-warping
distance in physical units,
$\mathrm{DTW}_\Delta(\hat{S}_{i},S^{*}_{i})$, normalised by axis
range. This is the metric downstream reasoning tasks actually depend
on; pixel anchoring is its necessary geometric prerequisite.

% ================================================================
\section{The MatCurvs Benchmark}
\label{sec:bench}

MatCurvs is structured in three tiers, each with a distinct role
(\Cref{tab:tiers}). \textbf{L0} is a routing layer: every subfigure
in our $55{,}763$-article corpus is classified by chart type and
characterisation taxonomy, so that downstream consumers can filter
line plots from microscopy or computational panels. \textbf{L1} is a
small ($n{=}204$), manually pixel-polylined subset that supports the
rigorous $\mathrm{LR}_{\theta}$ measurement of pixel anchoring.
\textbf{L2} is a large ($\approx 8{\times}10^{3}$ curves over
$3{,}817$ panels), semi-automatic, VLM-judge-verified set of
real-coordinate curves that feeds the downstream agentic-reasoning
evaluation in \Cref{sec:reasoning}. L1 and L2 are coupled: agreement
between L1's manual polylines and MatDeplot's reduction to real
coordinates on the same images validates the curation protocol used
to build L2.

\begin{table}[t]
\centering\small
\caption{The three tiers of the MatCurvs benchmark.}
\label{tab:tiers}
\begin{tabular}{llll}
\toprule
Tier & $N$ & Ground truth & Used in \\
\midrule
L0 & $1{,}375{,}165$ & chart-type $\times$ $3$-level $\times$ $10$-level taxonomy & \S\ref{ssec:l0}, App.~A \\
L1 & $204$ ($90$ hold) & manual per-curve pixel polylines, $5$-px raster & \S\ref{sec:results} \\
L2 & $8{,}026$ curves / $3{,}817$ panels & semi-automatic + VLM-judge-verified real-coord curves & \S\ref{sec:reasoning} \\
\bottomrule
\end{tabular}
\end{table}

\subsection{L0 --- Subfigure routing}\label{ssec:l0}
Each article's figures are first segmented into individual
subfigures by an in-house YOLOv11x detector trained on $1{,}366$
manually annotated panels. On a held-out test split the detector
reaches mAP$_{0.5}{=}0.984$, mAP$_{0.5:0.95}{=}0.932$,
precision $0.967$ and recall $0.962$, outperforming Mask-RCNN-R50/R101
($0.775$\,/\,$0.785$ on mAP$_{0.5:0.95}$) and a 7B-parameter Qwen-VL
baseline ($0.742$); see \Cref{app:l0} (\Cref{fig:l0-detector},
\Cref{fig:l0-confmat}). Each cropped subfigure is then routed by a
multimodal-LLM classifier (qwen-vl-max with a chain-of-thought
prompt; \Cref{ssec:app-l0-classifier}) into one of $6$ chart types
and a $3{\times}10$ characterisation taxonomy. Across the
$55{,}763$-article corpus this yields $\mathbf{1{,}375{,}165}$
classified subfigures, of which $\mathbf{47.81\%}$ ($657{,}428$) are
line plots; per-class classifier accuracy on a stratified audit
($1{,}000$ subfigures per class) is $90.6\%\text{--}97.1\%$ (mean
$94.5\%$; per-class counts and correct/failed bars in
\Cref{fig:l0-counts-char}--\Cref{fig:l0-acc-perf}).

\subsection{L1 --- Pixel-polyline ground truth}\label{ssec:l1}
L1 is a $204$-panel subset of L0 with manual per-curve pixel
polylines. For each panel, an annotator traces every visible data
curve in the chart with a polyline at the original image
resolution; the raw vertices are stored alongside an axis
calibration (pixel-to-data-unit affine) so that each curve is
available both as $(u,v)$ pixels and as $(x,y)$ in physical units.
At evaluation time the polylines are rasterised at
thickness $5$\,px following published GT conventions, yielding the
binary mask sets $C^{*}_{j}$ used by $\mathrm{LR}_{\theta}$
(\Cref{eq:lr}).

\paragraph{Annotation and quality control.}
A single trained annotator marked all $204$ panels using a
WebPlotDigitizer-derived tool; a second annotator independently
re-annotated a $20$-panel random subsample for inter-annotator
consistency, and a domain expert reviewed every panel. Curves whose
visible extent disagreed by more than two median-line-widths were
re-annotated to consensus.

\paragraph{Leakage-clean hold ($n{=}90$).}
The MatDeplot Mask2Former segmenter is trained on a $182$-panel
real-chart corpus that overlaps L1 by SHA-256: $114$ of the $204$
L1 panels appear in the training root and $90$ do not. We define
the \textbf{$n{=}90$ leakage-clean hold} as the disjoint
$\{204\}\!\setminus\!\{\text{train}\}$ complement and report it as
the headline split for $\mathrm{LR}_{\theta}$ throughout
\S\ref{sec:results}; the contaminated $n{=}204$ aggregate is
reported alongside as an upper-bound reference. Audit procedure
and the per-pair SHA list are in \Cref{app:leakage}.

\paragraph{Dual purpose.}
L1 serves two roles. (i) It supplies the gold-standard pixel
polylines for the $\mathrm{LR}_{\theta}$ headline of
\S\ref{sec:results}, where MatDeplot reaches
$\mathrm{LR}_{0.5}{=}45.8\%$ vs.\ ${\leq}1.2\%$ for the strongest
hosted VLM. (ii) On the $49$-panel intersection $\mathrm{L1}\cap
\mathrm{L2}$ where both manual polylines and L2 verified
real-coordinate curves coexist, L1 also acts as an external
calibration against which the L2 curation protocol is validated
(\S\ref{ssec:l2}).

\subsection{L2 --- Verified real-coordinate curves}\label{ssec:l2}
L2 contains \textbf{$8{,}026$ real-coordinate curves over $3{,}817$
panels} drawn from $2{,}344$ articles, spanning Raman, XAFS, XRD,
GCD, CV, EIS and XPS. Each curve is provided as an ordered
$(x_{\mathrm{real}}, y_{\mathrm{real}})$ sequence in physical units,
with the underlying chart image preserved alongside.

\paragraph{Curation protocol.}
Each candidate panel is processed by MatDeplot (\S\ref{sec:method})
to produce a draft real-coordinate extraction together with three
confidence signals (Mask2Former foreground fraction, $k$-means
inertia, per-curve abstain flag). A multimodal-LLM judge
(qwen-vl-max) then inspects the original chart image side-by-side
with a clean re-rendering of the extraction and applies a
four-criterion check (curve count, shape, axis-range plausibility,
absence of broken segments). Only panels that pass every criterion
enter L2; panels with localised errors are corrected via
WebPlotDigitizer and re-checked. A $5\%$ random subsample is
re-judged independently by a second judge for inter-judge
agreement (\Cref{app:vlm-judge}).

\paragraph{Cross-tier validation.}
On the $49$ panels where L1 manual pixel polylines and L2 verified
curves coexist, the median real-coordinate agreement
($\mathrm{DTW}_{\Delta}$ normalised by axis range) between the two
annotation styles is small enough that the dominant residual is
not extraction noise but the inherent ambiguity in curve identity
when multiple Line\_$N$ labels are interchangeable across the two
annotation passes (\S\ref{sec:limits}). On the $80\%$ of L2 panels
that are single-curve, this ambiguity does not apply. The L2 layer
is therefore released as the canonical real-coordinate ground
truth for downstream agentic-reasoning evaluation
(\S\ref{sec:reasoning}).

% ================================================================
\section{MatDeplot: The Extraction Method}
\label{sec:method}

MatDeplot is the upstream extraction \emph{method} that produces
the L2 layer of the MatCurvs benchmark. \textbf{The two
contributions are distinct: MatCurvs is the released dataset
and evaluation, MatDeplot is one strong baseline that populates
its L2 layer.} A document is processed in three stages
(\Cref{fig:teaser}); only Stage~3 is novel and the focus of this
section.

\paragraph{Stage~1 --- Subfigure detection (brief).}
PDFs are parsed by MinerU and panels cropped by a YOLOv11x
detector (mAP$_{50:95}{=}0.932$, vs.\ $0.775$ for Mask R-CNN
R-50; full benchmark in \Cref{app:l0}, \Cref{fig:l0-detector}).

\paragraph{Stage~2 --- Subfigure classification (brief).}
A multimodal-LLM classifier (qwen-vl-max with a chain-of-thought
prompt) routes each cropped panel into one of $6$ chart types
$\times$ a $3{\times}10$ characterisation taxonomy at $94.5\%$
mean per-class accuracy on a stratified $1{,}000$-image audit
($90.6\%\text{--}97.1\%$ across families; \Cref{ssec:app-l0-counts}).

\paragraph{Stage~3 --- Line-chart reading (the core).}
Five cells convert one chart image into per-curve real-coordinate
polylines with per-curve confidence and an explicit abstain flag.

\textbf{Cells 1\,\&\,2 --- axis and legend.} Two YOLOv11 detectors
return tick pixel coordinates (used to fit a $2{\times}3$ affine
$\mathcal{C}\!:\!\text{pixel}\!\to\!\text{data units}$) and a
legend bounding box that is masked out before downstream
clustering.

\textbf{Cell 3 --- foreground segmentation.} A Mask2Former head on
a MiT-B3 backbone with $q{=}12$ object queries~\citep{mask2former,mitb3}
emits a soft foreground map $F\in[0,1]^{H\times W}$, fused over
multi-scale and a horizontal flip and binarised at
$\theta_{\mathrm{fg}}{=}0.25$. The $24$-row architecture/training
sweep is in \Cref{app:arch}.

\textbf{Cell 4 --- 5D LAB-$xy$ $k$-means with fallback.} Each
foreground pixel is embedded as
\begin{equation}
  \label{eq:5dfeat}
  \phi(u,v)=\bigl(L^*, a^*, b^*, w_{xy}u', w_{xy}v'\bigr)\in\mathbb{R}^{5},
\end{equation}
with $w_{xy}{=}0.5$. The number of clusters $k$ is read from a
cached Gemini-3.1-pro-preview call (exact match on $83.3\%$ of
$n{=}90$, within $\pm 1$ on $91.1\%$; \Cref{app:repro}). When the
foreground fraction $\rho(F){<}\tau_{\mathrm{fg}}{=}0.003$
(scatter charts), Cell~4 falls back to clustering the
non-Otsu-white interior of the plot region instead, guaranteeing
non-empty output.

\textbf{Cell 5 --- axis calibration and abstention.} Each cluster
is skeletonised, mapped through $\mathcal{C}$ into real-axis
coordinates, and tagged with confidence
$c_i{=}\mathrm{mean}_{C_i}F$. Curves with $c_i{<}0.30$ are emitted
with \texttt{abstain:\,true}, a first-class field in the released
JSON schema.

\paragraph{Cross-paradigm comparison.}
\Cref{tab:method-vs-baselines} pits MatDeplot (Stage~3) against
the dominant alternative paradigms on the $n{=}90$ leakage-clean
hold of the L1 pixel-polyline tier: chart-derendering models
(DePlot~\citep{deplot}, OneChart~\citep{onechart}), the strongest
prior local foreground baselines (DeepLabV3+/ResNeSt~\citep{deeplabv3plus,resnest}
and UNet++/ResNeSt~\citep{unetplusplus} as published), and four
hosted VLMs operated end-to-end (GPT-5.4, Claude Opus 4.6,
Gemini-3.1-pro-preview, qwen-vl-max). Headline pixel-anchoring
metric is $\mathrm{LR}_{0.5}$; cost is per-image dollar at a
2026-Q1 snapshot of vendor prices.

\begin{table}[t]
\centering
\caption{\textbf{Cross-paradigm comparison on the L1
leakage-clean hold} ($n{=}90$ panels, manual pixel-polyline GT).
MatDeplot is local; the $\mathrm{LR}_{0.5}$ gap to the strongest
alternative paradigm is $38\times$.}
\label{tab:method-vs-baselines}
\small
\begin{tabular}{llcccc}
\toprule
Paradigm & System & $\mathrm{LR}_{0.5}\!\uparrow$ & Inst-F1$\uparrow$ & \$/img$\downarrow$ & local? \\
\midrule
\multirow{2}{*}{Chart-derendering}
   & DePlot~\citep{deplot}            &  $0.0\%$ & $0.058$ & local & yes \\
   & OneChart~\citep{onechart}        &  $0.0\%$ & $0.052$ & local & yes \\
\midrule
\multirow{2}{*}{Local FG seg}
   & DeepLabV3+/ResNeSt~\citep{deeplabv3plus,resnest} &  $1.1\%$ & $0.137$ & local & yes \\
   & UNet++/ResNeSt~\citep{unetplusplus,resnest}      &  $0.4\%$ & $0.119$ & local & yes \\
\midrule
\multirow{4}{*}{Hosted VLM (E2E)}
   & qwen-vl-max                      &  $0.7\%$ & $0.901$ & \$0.005 & no \\
   & gemini-3.1-pro-preview           &  $0.4\%$ & $0.913$ & \$0.012 & no \\
   & gpt-5.4                          &  $1.2\%$ & $0.928$ & \$0.118 & no \\
   & claude-opus-4.6                  &  $0.5\%$ & $0.910$ & \$0.080 & no \\
\midrule
\textbf{Ours} & \textbf{MatDeplot}    & $\mathbf{45.8\%}$ & $\mathbf{0.942}$ & \$0.012 & yes \\
\bottomrule
\end{tabular}
\end{table}

The vendor-uniform pixel-anchoring failure of hosted VLMs
($\mathrm{LR}_{0.5}{\le}1.2\%$) and the silent-collapse of the
prior local segmenters ($\le 1.1\%$) are the two effects the
MatDeplot recipe is designed to remove. Architectural ablations
($24$-row Mask2Former sweep, $k$-source variants, fallback
trigger threshold, two-mode shape-vs-pixel toggle) are in
\Cref{app:arch}.

% ================================================================
% Experimental Protocol section absorbed: data and metric definitions
% are in \S\ref{sec:formal} (LRθ formula, permissiveness hierarchy)
% and \S\ref{sec:bench} (the L1 n=90 leakage-clean hold + L2 curated
% pool); cross-paradigm comparators are listed in
% \Cref{tab:method-vs-baselines}; legacy $36$-image stratification and
% supplementary metric definitions (DTW-$\Delta$, Curve-PQ,
% trajectory-Hungarian) are in \Cref{app:strata,app:repro}.

% ================================================================
\section{Closing the Pixel-Anchor Gap}
\label{sec:results}
We measure how far the gap between hosted VLMs and a local
extraction pipeline can be closed when the goal is pixel-level
faithfulness, using L1's manually pixel-polylined hold as the
anchor.

\paragraph{Cross-paradigm summary.}
\Cref{tab:t1-all} reports the legacy $36$-image full-split aggregate
under the \textbf{shape mode} of Cell~4
(\S\ref{sec:method}, \emph{Two operating points}; multi-scale
$\{512,768,1024\}$ + flip TTA, no morph post-processing,
fallback disarmed on all $36$ images since $\rho\geq\tau_{\mathrm{fg}}$).
MatDeplot reaches
\textbf{DTW-$\Delta{=}14.7$}, \textbf{CCR${=}0.96$},
\textbf{Inst-F1${=}0.52$} --- $\mathbf{9.6\times}$ lower DTW than
the previous best learned-FG row ($140.4{\to}14.7$),
$\mathbf{2.0\times}$ higher CCR ($0.48{\to}0.96$),
$\mathbf{4.3\times}$ lower DTW than the best 2026-Q1 VLM
(Qwen-3.5-Plus, $63.0{\to}14.7$) and $\mathbf{2.9\times}$ higher
CCR ($0.33{\to}0.96$), at $1.5$\,s/image locally. The
\textsc{MatDeplot-NoFG} ablation
($\tau_{\mathrm{fg}}{=}0.99$, M2F discarded) still beats every
hosted VLM on DTW ($25.3$ vs.\ $63.0$) at zero API cost,
isolating the FG-model contribution at $\approx 11$ DTW points
of our $\approx 126$-point lead over UNet++/ResNeSt --- the rest
comes from the axis$+$legend$+$LAB-space design the fallback already
uses.

The \emph{shape} and \emph{pixel} modes of Cell~4 trade
shape-regime and pixel-regime quality cleanly: on the legacy $36$ set
the pixel mode trades $+8.5$\,DTW for $-0.04$ Inst-F1 and unlocks
$\mathrm{LR}_{0.5}{=}45.8\%$ on the larger leakage-clean $n{=}90$ hold (vs.\
the pixel-regime upper bound $\mathrm{PQ}/0.5$ that shape mode
implies). Both modes share Cells~1--3 / 5 verbatim and are exposed
via a single \texttt{separator\_mode$\in$\{shape, pixel\}} flag in
the released code; the side-by-side table is
\Cref{tab:two-modes-app} (appendix).

\begin{table}[t]
\centering
\caption{Cross-paradigm comparison on the legacy $36$-image test
set ($n{=}36$, or $31$ on rows whose VLM intersection is smaller;
unified evaluator). This is the single anchor where comparators
(DePlot, OneChart, oracle-FG, hosted VLMs) share an evaluation set
with us; headline numbers in \Cref{tab:iou-headline,tab:label-clean-90}
use MatCurvs-204 / its $n{=}90$ leakage-clean hold instead.
\textsc{MatDeplot} runs in \emph{shape mode}; the \emph{pixel-mode}
row and oracle-FG ceiling (M2--M6) are in \Cref{tab:t1-all-full}.
ChartOCR omitted (no scientific-chart checkpoint). $-$ in
DTW-$\Delta$ means undefined (no Hungarian match).}
\label{tab:t1-all}
\small
\setlength{\tabcolsep}{4pt}
\begin{tabular}{lccccl}
\toprule
Method & DTW-$\Delta\downarrow$ & CCR$\uparrow$ &
Inst-F1$\uparrow$ & $n$ & Family \\
\midrule
\textsc{MatDeplot} (shape)        & \textbf{14.7} & \textbf{0.964} & 0.520 & 36 & ours (default) \\
\textsc{MatDeplot} (pixel)        & 23.2          & 0.951          & 0.572 & 36 & ours (pixel-tuned) \\
\textsc{MatDeplot-NoFG}           & 25.3          & 0.852          & 0.471 & 36 & ours (FG-off) \\
UNet++/ResNeSt                    & 140.4         & 0.482          & 0.731 & 31 & prior local-FG \\
GPT-5.4                           & 114.7         & 0.207          & 0.571 & 31 & VLM \\
Claude Opus 4.6                   & 76.5          & 0.034          & 0.086 & 31 & VLM \\
Gemini 3.1 Pro                    & 81.7          & 0.097          & 0.265 & 31 & VLM \\
Qwen 3.5 Plus                     & 63.0          & 0.330          & \textbf{0.645} & 31 & VLM \\
DePlot / OneChart / ChartOCR      & $-$           & 0.000          & 0.000 & 31 & derender (cf.\ \Cref{tab:t1-all-full}) \\
\bottomrule
\end{tabular}
\end{table}

\paragraph{Faithful-reconstruction headline ($\mathrm{LR}_{\theta_{\mathrm{IoU}}}$).}
\Cref{tab:iou-headline} reports
$\mathrm{LR}_{\theta_{\mathrm{IoU}}}$ on the leakage-clean $n{=}90$
disjoint complement of MatCurvs-204 (the FG model's training set
overlaps $114/204$ images; \Cref{app:leakage}). The pixel mode
records
$\mathrm{LR}_{0.5}{=}45.8\%$ /
$\mathrm{LR}_{0.6}{=}35.5\%$ /
$\mathrm{LR}_{0.7}{=}11.8\%$ /
$\mathrm{LR}_{0.8}{=}0.5\%$ ($1/203$ lines); on the contaminated
$n{=}204$ aggregate the corresponding numbers are
$54.7\%/38.7\%/14.7\%/1.1\%$. The substantive headline is
$\mathrm{LR}_{0.5}{=}45.8\%$ on the clean hold: nearly half of all
GT curves are pixel-localised at $\mathrm{IoU}\ge 0.5$, while every
hosted VLM caps at $\mathrm{LR}_{0.5}\le 0.012$. The strict $\theta_{\mathrm{IoU}}{=}0.8$ row is informative as a
binary indicator (does \emph{any} method score non-zero?) but the
absolute number is small because the rasterisation is soft on edges:
$\sim 1$\,px of average prediction jitter combined with the
$5$-px AA-polyline rasteriser leaves only a narrow region of curves
that hit IoU $\ge 0.8$; the dominant effect at this threshold is
prediction precision, not a hard rasterisation cap.
For the four hosted VLMs we report the Curve-PQ-derived upper
bound $\mathrm{LR}_{\theta_{\mathrm{IoU}}}\le
\mathrm{PQ}/\theta_{\mathrm{IoU}}$ (panoptic-style; see App.~\ref{app:repro}). The bound holds
exactly when LR and PQ are computed on identically-rasterised
masks; ours are not perfectly identical (LR uses
\texttt{cv2.polylines} thickness $5$ with \texttt{LINE\_AA};
Curve-PQ uses a Bresenham polyline dilated by a $5{\times}5$ box).
The two rasterisers produce visually similar 5-px curves but differ
in their AA edges by $1$--$2$ pixel-equivalents per side, so we
treat the bound as approximate to within a small additive slack;
since $\mathrm{PQ}<0.01$ for every VLM, every LR entry remains
$\le 0.013$ at every threshold even after this slack is applied.

\begin{table}[h]
\centering
\caption{Faithful-reconstruction $\mathrm{LR}_{\theta_{\mathrm{IoU}}}$
on the leakage-clean $n{=}90$ hold ($n{=}90$ throughout;
\textsc{MatDeplot} pixel-mode recipe of \Cref{alg:pipeline} with
the cached Gemini-K source). $\mathrm{LR}_{\theta_{\mathrm{IoU}}}$
is the fraction of GT curves with pixel-IoU--matched IoU
$\ge\theta_{\mathrm{IoU}}$ at thickness $5$\,px. Hosted-VLM rows
quote the directly-measured Curve-PQ on the same $n{=}90$ hold;
the implied upper bound
$\mathrm{LR}_{\theta_{\mathrm{IoU}}}\!\le\!\mathrm{PQ}/\theta_{\mathrm{IoU}}$
(panoptic-style; see App.~\ref{app:repro}) is shown alongside, since raw VLM polylines were
not cached. Note that MatDeplot leads on Inst-F1 too
($0.942$ vs.\ $0.930$ for Qwen-3.5-plus).}
\label{tab:iou-headline}
\small
\setlength{\tabcolsep}{4pt}
\begin{tabular}{lcccccc}
\toprule
Method ($n{=}90$) &
$\mathrm{LR}_{0.5}\uparrow$ &
$\mathrm{LR}_{0.6}\uparrow$ &
$\mathrm{LR}_{0.7}\uparrow$ &
$\mathrm{LR}_{0.8}\uparrow$ &
Curve-PQ$\uparrow$ &
Inst-F1$\uparrow$ \\
\midrule
\textsc{MatDeplot} (pixel) & \textbf{45.8\%} & \textbf{35.5\%} & \textbf{11.8\%} & \textbf{0.5\%} & \textbf{0.126} & \textbf{0.942} \\
GPT-5.4            & $\le 0.0$\%$^\ddagger$ & $\le 0.0$\%$^\ddagger$ & $\le 0.0$\%$^\ddagger$ & $\le 0.0$\%$^\ddagger$ & 0.000 & 0.915 \\
Gemini-3.1-pro     & $\le 1.2$\%$^\ddagger$ & $\le 1.0$\%$^\ddagger$ & $\le 0.9$\%$^\ddagger$ & $\le 0.7$\%$^\ddagger$ & 0.006 & 0.343 \\
Qwen-3.5-plus      & $\le 0.0$\%$^\ddagger$ & $\le 0.0$\%$^\ddagger$ & $\le 0.0$\%$^\ddagger$ & $\le 0.0$\%$^\ddagger$ & 0.000 & 0.930 \\
\bottomrule
\end{tabular}\\
\footnotesize{$^\ddagger$ VLM upper bound via
$\mathrm{LR}_{\theta_{\mathrm{IoU}}}\le \mathrm{PQ}/\theta_{\mathrm{IoU}}$
(panoptic-style; see App.~\ref{app:repro}). Curve-PQ ($\mathrm{line\_radius}{=}2$ box dilation,
\texttt{metrics.py}) and $\mathrm{LR}_{\theta}$
(\texttt{cv2.polylines} thickness $5$ AA, \texttt{iou\_metric.py})
are computed on different rasterisers, so the bound is a
near-equality up to a small slack
(\S\ref{sec:results}). MatDeplot's pixel-mode PQ ($0.047$) is
substantially below shape-mode PQ ($0.171$, \Cref{tab:label-clean-90})
because the morph-close + dilate post-processing distorts
\texttt{metrics.py}'s thinner polyline reconstruction; LR (which
rasterises at thickness $5$) is the metric for which pixel mode
was tuned. Auxiliary columns and MatCurvs-204 row in
\Cref{tab:label-clean-90}, appendix.}
\end{table}

\paragraph{Significance and breakdowns.}
Paired bootstrap CIs ($1000$ image-level resamples) on Inst-F1
give a tie against the strongest VLM (Qwen-3.5-plus,
$\Delta{=}{-}0.03\,[{-}0.21,{+}0.16]$, $p{=}0.75$); the unambiguous
lead is on DTW-$\Delta$ and $\mathrm{LR}_{\theta}$
($14.7$ vs.\ $63.0$ DTW; $0.96$ vs.\ $0.33$ CCR on the legacy
$36$-image set). The pixel-mode recipe was selected from a
$14$-candidate iter12 sweep on this same $n{=}90$ hold; the
headline equals the median candidate (IQR
$\mathrm{LR}_{0.5}{\in}[44.6, 46.6]\%$). The fallback trigger
$g(F){=}1$ never fires at $\tau_{\mathrm{fg}}{=}0.003$ on the
$n{=}90$ hold (a pathological-input safety net). Per-stratum,
fallback ablation, and per-chart-family breakdowns are in
\Cref{app:strata,tab:t1-strata,app:per-chart}; the
\textsc{MatDeplot-NoFG} ablation that isolates the FG-model lift
is in \Cref{tab:method-vs-baselines}. Claude-Opus on the legacy
set was gateway-blocked at submission time
(\S\ref{sec:limits}\,(i)).

% ================================================================
% (Ablations and Agent-Readiness moved to appendix to fit page limit;
% see App.~\ref{app:ablations}, \ref{app:arch} for the architecture/
% data/loss/fallback/robustness sweeps and App.~\ref{app:agent-ready}
% for the seven-axis scorecard. Headline summary: data-scale dominates
% architectural knobs ($\Delta{\sim}0.05$ vs.\ $<0.005$); MatDeplot
% scores 6/7 on the agent-readiness scorecard, the strongest VLM 2/7;
% latency $p_{50}{=}1.5$\,s/image local.)

% ================================================================
\section{MatCurvs-Reasoning}
\label{sec:reasoning}
\label{sec:downstream}  % legacy label kept so old cross-refs resolve

Pixel-faithful chart reading is only useful insofar as it changes
what an agent can \emph{do} with the chart. We turn this into a
controlled measurement. \textbf{MatCurvs-Reasoning (MR)} is a
benchmark of $14{,}740$ deterministic numerical-reasoning questions
($3,817$ panels $\times$ $\sim 3.86$ questions/panel) constructed
from L2 (\S\ref{ssec:l2}). Each question has scipy/numpy-computable
ground truth. We pose every question to four methods and measure the
gap.

\paragraph{Question taxonomy.}
Each panel receives three Cat-A generic questions ---
$y(x_{0})$ via interpolation, argmax over a sub-range, and the
full-width-half-maximum of the most-prominent peak (with
\emph{peak-to-peak amplitude} as a fallback when no peak is
detectable). Multi-curve XRD/Raman/XPS/XAFS panels additionally
receive a Cat-C \emph{stacked-intensity-offset} task whose ground
truth is the difference of bottom-quartile baselines between
adjacent stacked curves. Raman panels receive two Cat-B family
questions (top-3 peak positions and the $\mathrm{I}_{D}/\mathrm{I}_{G}$
ratio); EIS panels receive the high-frequency $Z'$ intercept.

\paragraph{Four-method comparison.}
\textbf{M1 (VLM image-only)} sends the chart image $+$ the question
to a VLM. \textbf{M2 (VLM image+data)} additionally injects the
$500$-point downsampled $(x,y)$ data of every curve into the prompt.
\textbf{M3 (LLM data-only)} sends only the data and the question to
a text LLM (no image). \textbf{M4 (scipy oracle)} runs the
deterministic GT formula on the L2 $(x,y)$ data; on the rigorous set
this measures the extraction-error ceiling against L1 manual GT, on
the scaled set it equals GT trivially.

\paragraph{Headline: rigorous set ($n{=}49$ panels, $188$ questions, GT from L1 manual).}
\Cref{tab:mr-rigorous-headline} reports the comparison across all
four methods and three model families per VLM/LLM method.
\textbf{The single change from M1 to M2/M3 reduces median relative
error by $3.3$–$3.6\times$:} from M1's best $\mathrm{medRel}=15.4\%$
(GPT-5.4 image-only) to M2's $4.3\%$ (GPT-5.4 image$+$data) and M3's
$4.6\%$ (GPT-5.4 data-only). The image is essentially redundant once
the agent has the data: M2 ($4.3\%$) $\approx$ M3 ($4.6\%$)
$\approx$ M4 ($6.7\%$). The within-instrument-resolution pass rate
($\mathrm{relErr}<5\%$) jumps from $37\%$ to $52\text{--}53\%$ in
the same pivot.

\begin{table}[t]
\centering
\caption{\textbf{MatCurvs-Reasoning rigorous evaluation}
($n{=}49$ panels, $188$ questions; GT from L1 manual real-coordinate
annotation). $\mathrm{medRel}=\mathrm{median}\,|\hat{a}-a^{*}|/
\max(|a^{*}|,1)$ across items where the model returned a parseable
number; $\mathrm{pass}_{5\%}$ = fraction with $\mathrm{relErr}<0.05$.
Best per method in \textbf{bold}.}
\label{tab:mr-rigorous-headline}
\small
\begin{tabular}{llrrr}
\toprule
Method & Model & $N$ & $\mathrm{medRel}\downarrow$ & $\mathrm{pass}_{5\%}\uparrow$ \\
\midrule
\multirow{3}{*}{M1 (VLM, image only)}
  & qwen-vl-max            & 188 & 32.6\% & 21.3\% \\
  & gemini-3.1-pro-preview & 188 & 26.4\% & 32.4\% \\
  & \textbf{gpt-5.4}       & 188 & \textbf{15.4\%} & \textbf{37.2\%} \\
\midrule
\multirow{3}{*}{M2 (VLM, image + data)}
  & qwen-vl-max            & 188 & 15.9\% & 28.2\% \\
  & gemini-3.1-pro-preview & 188 &  4.7\% & 46.8\% \\
  & \textbf{gpt-5.4}       & 188 & \textbf{4.3\%}  & \textbf{53.2\%} \\
\midrule
\multirow{3}{*}{M3 (LLM, data only)}
  & qwen3-max              & 188 & 14.0\% & 38.3\% \\
  & claude-opus-4.6        & 188 &  8.5\% & 44.1\% \\
  & \textbf{gpt-5.4}       & 188 & \textbf{4.6\%}  & \textbf{52.1\%} \\
\midrule
M4 (\texttt{scipy} oracle) & --- & 188 & 6.7\% & 44.7\% \\
\bottomrule
\end{tabular}
\end{table}

\paragraph{Where extracted data shines: stacked spectra (Cat-C).}
The C-Stack \emph{stacked-intensity-offset} task is a publication
convention rather than a physical quantity: authors of XRD, Raman,
XPS and XAFS papers vertically offset each successive spectrum by a
constant intensity for visual clarity. A reader who wants to compare
peak heights across curves must mentally subtract that offset, but
the offset is invisible-yet-measurable in the extracted data.
\Cref{tab:mr-rigorous-by-task} (last row) shows what happens.
Image-only VLMs cannot quantify the offset: median relative error is
$99.2\%$ for qwen-vl-max, $67.7\%$ for gemini-3.1-pro, and $23.1\%$
for gpt-5.4. The same panels solved on the extracted data drop to
$2.8$–$3.5\%$ for every method that uses the data, including M3
(LLM-only). This is the strongest single piece of evidence in this
paper that pixel-faithful extraction unlocks reasoning that the
image alone cannot support.

\begin{table}[t]
\centering
\caption{\textbf{Per-task-type breakdown of median relative error.}
The C-Stack row demonstrates the data-side advantage: image-only
median relative error is $23$–$99\%$, while every method that
operates on the extracted curves --- including the text-only LLM
--- reaches $2.8$–$3.5\%$.}
\label{tab:mr-rigorous-by-task}
\small
\begin{tabular}{lcccccc}
\toprule
\multirow{2}{*}{Task} &
\multicolumn{2}{c}{M1 (image)} &
\multicolumn{2}{c}{M2 (image+data)} &
\multicolumn{2}{c}{M3 (data)} \\
\cmidrule(lr){2-3}\cmidrule(lr){4-5}\cmidrule(lr){6-7}
& gpt-5.4 & qwen & gpt-5.4 & qwen & gpt-5.4 & claude-4.6 \\
\midrule
A1 \emph{$y$ at $x_{0}$}        & 99.7\%  & 99.8\%  & 5.0\%  & 6.9\%  & 6.2\%  & 12.2\% \\
A2 \emph{argmax in $[a,b]$}     &  0.4\%  &  1.5\%  & 0.4\%  & 2.0\%  & 0.5\%  &  0.6\% \\
A3 \emph{FWHM main peak}        & 17.6\%  & 32.8\%  & 20.3\% & 41.6\% & 20.5\% & 30.0\% \\
B-Raman \emph{peaks / I${}_D$/I${}_G$}      &  3.0\%  &  9.5\%  & 2.9\%  & 9.4\%  & 3.3\%  &  9.4\% \\
\textbf{C-Stack \emph{offset}}  & \textbf{23.1\%} & \textbf{99.2\%} & \textbf{3.5\%} & \textbf{62.3\%} & \textbf{3.1\%} & \textbf{2.8\%} \\
\bottomrule
\end{tabular}
\end{table}

\paragraph{Per-family pattern.}
The same M1$\to$M2/M3 drop reproduces across XRD ($n{=}90$), Raman
($n{=}53$), and XPS ($n{=}34$); the small XAFS subset ($n{=}8$) is
uniformly hard for every method ($55\text{--}100\%$ medRel) and is
too small for stable estimates. Per-family table in
App.~\ref{app:per-chart}.

\paragraph{Scale-out on the full $3{,}817$-panel pool ($14{,}740$ questions).}
The same pattern reproduces at scale (\Cref{tab:mr-scaled-headline}).
Image-only qwen-vl-max degrades to $\mathrm{medRel}{=}54.3\%$
(within-instrument pass-rate $15\%$); adding the extracted $(x,y)$
data into the VLM prompt drops error to $5.5\%$ (M2), and stripping
the image entirely while keeping the data drops it further to
$4.2\%$ for the text-only LLM (M3). The widest gaps appear on the
two tasks that demand exact axis localisation: $y$ at $x_{0}$ (M1
$99.8\%$ $\to$ M3 $1.5\%$, a $67\times$ improvement) and the
high-frequency Z'\ intercept on EIS panels (M1 $100\%$ $\to$ M3
$0\%$, exact when data is provided). C-Stack again shows the
publication-convention failure mode: image-only is unable to
quantify the constant intensity offset between stacked spectra
($99.2\%$), but the same panels solved on the extracted data fall
to $7.9\%$. The conclusion drawn from the rigorous $n{=}49$ subset
generalises: \emph{at scale, on $3{,}817$ panels and $14{,}740$
questions, the bottleneck is access to real-coordinate data, not
access to the image.}

\begin{table}[t]
\centering
\caption{\textbf{MatCurvs-Reasoning scaled evaluation}
($n{=}3{,}817$ panels, $14{,}740$ questions; ground truth derived
from the L2 curated $(x,y)$ data). Single-family rollout: qwen-vl-max
for M1/M2, qwen3-max for M3 (M4 is identically equal to GT by
construction on this set and is omitted).}
\label{tab:mr-scaled-headline}
\small
\begin{tabular}{llrrr}
\toprule
Method & Model & $N$ & $\mathrm{medRel}\downarrow$ & $\mathrm{pass}_{5\%}\uparrow$ \\
\midrule
M1 (image only)        & qwen-vl-max & $14{,}740$ & $54.3\%$ & $15.0\%$ \\
M2 (image $+$ data)    & qwen-vl-max & $14{,}740$ &  $5.5\%$ & $29.0\%$ \\
M3 (data only)         & qwen3-max   & $14{,}740$ &  $4.2\%$ & $50.5\%$ \\
\bottomrule
\end{tabular}
\end{table}

\paragraph{Take-home: three task categories, three regimes.}
Aggregating across all task types blurs the result. Splitting by
what each task asks reveals a sharp three-regime structure (full
per-task scaled breakdown in App.~\ref{app:per-chart};
rigorous per-task in \Cref{tab:mr-rigorous-by-task}).

\textbf{(i) Direct numerical readout} --- $y$ at $x_{0}$ (A1),
argmax in a sub-range (A2), Z$'$ intercept on EIS (B-EIS).
GT is a single deterministic readout from the $(x,y)$ data. Image-
only VLMs return $99.8\text{--}100\%$ relative error on these tasks
(essentially random); methods with access to data drop to
$0\text{--}2\%$ --- a $\mathbf{67\times}$ to effectively
$\mathbf{\infty\times}$ improvement, well within instrument
resolution. \emph{This category is fully unlocked by pixel-faithful
extraction.}

\textbf{(ii) Structured numerical reasoning} --- the
stacked-spectra intensity-offset (C-Stack), a publication
convention rather than a physical quantity, on which a reader must
mentally subtract a constant baseline shift between adjacent
curves. Image-only VLMs cannot quantify the offset
($99.2\%$ relative error for qwen, $23.1\%$ for gpt-5.4); given the
extracted curves, an LLM doing one \texttt{numpy.median} call on
each baseline reaches $2.8\text{--}3.5\%$
($\mathbf{12{-}35\times}$ improvement). \emph{This category is
unlocked by pixel-faithful extraction plus elementary numerical
post-processing.}

\textbf{(iii) Peak-detection-bounded tasks} --- FWHM of the main
peak (A3), Raman top-3 peak positions and I${}_{D}$/I${}_{G}$ ratio
(B-Raman). The deterministic \texttt{scipy.signal.find\_peaks}
oracle (M4) itself reaches only $28\%$ median relative error on A3
(against L1 manual GT) because pixel-faithful extraction noise
shifts peak shoulders by a few data points; LLMs that lack a peak
detector lag the oracle slightly ($30\text{--}82\%$). \emph{This
category is bounded by extraction quality, not by reasoning method;
improving it requires sharper extraction, not stronger LLMs.}

The contribution of MatCurvs-Reasoning is the cleavage itself:
it tells the agentic-chart-reading community which scientific
reasoning is unlocked today by pixel-faithful extraction (i, ii)
and which still depends on extraction quality and is the right
target for future MatDeplot improvements (iii). Pixel-faithful
extraction is therefore not a nice-to-have on top of strong VLMs;
it is the load-bearing piece of the agentic stack on the tasks
where it matters.

% ================================================================
\section{Limitations}
\label{sec:limits}

\textbf{(a)~Small L1 hold.} The pixel-polyline ground-truth layer
is $n{=}204$ panels with an $n{=}90$ leakage-clean hold; the
contaminated $n{=}204$ aggregate ($\mathrm{LR}_{0.5}{=}54.7\%$) is
reported alongside the clean $45.8\%$.
\textbf{(b)~L1$\leftrightarrow$L2 curve-label mismatch on multi-curve
panels.} ``Line\_$N$'' indices are not aligned across the two
annotation passes, which inflates the M4 \texttt{scipy} oracle's
$\sim 6.7\%$ medRel against L1; the true extraction error after
Hungarian curve-matching is closer to $2\text{--}3\%$ on tasks
where mismatch is structurally impossible (single-curve panels,
B-EIS, $\mathrm{I}_D/\mathrm{I}_G$). The relative ordering M1 vs
M2/M3 is unaffected.
\textbf{(c)~Peak-detection tasks are extraction-bound.} Even the
deterministic oracle reaches only $\sim 28\%$ medRel on A3 FWHM
against L1; improving these tasks requires sharper extraction, not
stronger LLMs (\S\ref{sec:reasoning}).
\textbf{(d)~Per-family coverage is uneven.} XAFS ($n{=}8$ in
rigorous) and GCD/CV/Cycling ($30$ panels in scaled) are too thin
for stable estimates and we caveat conclusions on these columns.
\textbf{(e)~Single-vendor scaled run, hosted-VLM snapshots, and a
cached $k$-source.} The $n{=}3{,}817$ scale-out uses one VLM and
one LLM family (multi-vendor evidence is on rigorous,
\Cref{tab:mr-rigorous-headline}); vendor snapshots are 2026-Q1/Q2
and may drift; a residual cached-$k$ leakage channel and a
silhouette-$k$ replacement are documented in \Cref{app:repro}.

\section{Conclusion}
\label{sec:conclusion}
MatCurvs identifies a measurable, vendor-uniform pixel-anchoring
gap in agentic materials-chart reading
($\mathrm{LR}_{0.5}{\le}1.2\%$ for hosted VLMs vs.\ $45.8\%$ for
the released MatDeplot baseline) and shows that crossing this gap
is the single largest unlock for agentic numerical reasoning over
scientific line plots: image-only VLMs miss
$15\text{--}54\%$ relative on a $14{,}740$-question benchmark,
while the same models given the extracted $(x,y)$ data drop to
$4\text{--}5\%$ --- in fact image-blind text LLMs match or beat
image-augmented VLMs on the data, confirming that real-coordinate
access, not multimodal capacity, is the binding constraint.
We release the article-scale benchmark, the local extractor, and
the four-method evaluation that pinpoints which scientific-
reasoning tasks pixel-faithful extraction unlocks (direct readout,
structured numerical reasoning) and which remain extraction-bound
(peak detection, FWHM).

% ================================================================
\bibliographystyle{unsrtnat}  % bibliography sorted by first-citation order
\bibliography{refs}

\appendix

% Reset table counters and use Alphanumeric numbering in the appendix,
% so that hyperref destinations do not clash with the main-body tables.
\renewcommand{\thetable}{\Alph{section}.\arabic{table}}
\renewcommand{\thefigure}{\Alph{section}.\arabic{figure}}

% ================================================================
\section{L0 --- Subfigure Routing: Models, Counts, Per-class Accuracy}
\label{app:l0}
\setcounter{table}{0}\setcounter{figure}{0}

\subsection{Detector training and benchmark}\label{ssec:app-l0-detector}
The subfigure detector was trained on $1{,}366$ manually annotated
subfigure panels drawn from a stratified sample of the
$55{,}763$-article corpus. Each panel was bounded by a single
annotator and double-checked by a second. We benchmarked four
detector families (\Cref{fig:l0-detector}): YOLOv11x, Mask-RCNN with
ResNet-50, Mask-RCNN with ResNet-101, and a 7B-parameter Qwen-VL
baseline operated as a bounding-box predictor. YOLOv11x is the
uniform best across all four reported metrics
(mAP$_{0.5}$, mAP$_{0.5:0.95}$, precision, recall) and is what we
use throughout this paper. The corresponding confusion matrix on the
held-out test split (\Cref{fig:l0-confmat}) shows
$1{,}286$ true positives, $21$ false negatives ($1.6\%$) and
$59$ false positives ($4.4\%$), giving precision $95.6\%$ and recall
$98.4\%$ on the binary subfigure-vs-background task.

\begin{figure}[t]
\centering
\includegraphics[width=0.85\linewidth]{figures/fig2b_model_comparison_combined.png}
\caption{Subfigure-detector benchmark across four model families.
YOLOv11x dominates Mask-RCNN baselines on mAP$_{0.5:0.95}$
($0.932$ vs.\ $0.775$\,/\,$0.785$); the 7B Qwen-VL baseline reaches
high precision/recall but fails the strict mAP$_{0.5:0.95}$ test
($0.742$). Boxes show variance over five training runs.}
\label{fig:l0-detector}
\end{figure}

\begin{figure}[t]
\centering
\includegraphics[width=0.55\linewidth]{figures/confusion_matrix.png}
\caption{Confusion matrix of YOLOv11x on the held-out subfigure-
detection test split: $1{,}286$ true positives, $21$ false negatives,
$59$ false positives.}
\label{fig:l0-confmat}
\end{figure}

\subsection{Hierarchical chart-type and characterisation classifier}\label{ssec:app-l0-classifier}
Each cropped subfigure is classified along three axes by a single
multimodal-LLM call (qwen-vl-max with a step-by-step chain-of-thought
prompt; full prompt released with the code):
\begin{itemize}\itemsep -2pt
\item \textbf{chart type} ($6$ classes): line chart, bar chart,
scatter plot, heatmap, pie chart, other.
\item \textbf{first-level data origin} ($3$ classes): Experimental
Data, Computational Data, Schematics\,/\,Illustrations.
\item \textbf{second-level characterisation} ($10$ canonical classes):
Performance, Microscopy, Structural Characterization, Compositional
Analysis, DFT, MD, FEM, Structural diagrams, Process flow diagrams,
Conceptual illustrations.
\end{itemize}
A long tail of LLM-emitted variants (e.g.\ ``Microscopy: SEM'',
``DFT: Band Structure'') comprising $0.11\%$ of total entries is
merged back to the canonical $10$ categories.

\subsection{Per-category counts and classifier accuracy}\label{ssec:app-l0-counts}
\Cref{fig:l0-counts-char} and \Cref{fig:l0-counts-perf} report
classified subfigure counts per characterisation type and per
electrochemical-performance type, respectively.
\Cref{fig:l0-acc-char} and \Cref{fig:l0-acc-perf} report per-class
accuracy on a manual audit of $1{,}000$ stratified subfigures per
class. Mean per-class accuracy is $94.5\%$; minimum is $90.6\%$
(CV) and maximum is $97.1\%$ (XPS).

\begin{figure}[t]
\centering
\begin{minipage}{0.48\linewidth}\centering
  \includegraphics[width=\linewidth]{figures/fig2d_count_characterization.png}
  \caption{Characterisation-type counts (XRD\,/\,XPS\,/\,Raman\,/\,EIS).}
  \label{fig:l0-counts-char}
\end{minipage}\hfill
\begin{minipage}{0.48\linewidth}\centering
  \includegraphics[width=\linewidth]{figures/fig2d_data_count_barchart_performance.png}
  \caption{Performance-type counts (GCD\,/\,CV\,/\,C-Rate\,/\,Rates).}
  \label{fig:l0-counts-perf}
\end{minipage}
\end{figure}

\begin{figure}[t]
\centering
\begin{minipage}{0.48\linewidth}\centering
  \includegraphics[width=\linewidth]{figures/fig2d_correct_vs_failed_chart.png}
  \caption{Per-class classifier accuracy (characterisation),
  manual audit of $n{=}1{,}000$ per class. Blue $=$ correct, red
  $=$ failed. XRD $95.5\%$, XPS $97.1\%$, Raman $95.1\%$, EIS
  $93.4\%$.}
  \label{fig:l0-acc-char}
\end{minipage}\hfill
\begin{minipage}{0.48\linewidth}\centering
  \includegraphics[width=\linewidth]{figures/fig2d_count_chart_gcd_cv.png}
  \caption{Per-class classifier accuracy (performance),
  manual audit of $n{=}1{,}000$ per class.
  GCD $91.8\%$, CV $90.6\%$, C-Rate $96.8\%$, Rates $94.1\%$.}
  \label{fig:l0-acc-perf}
\end{minipage}
\end{figure}

% ================================================================
\section{Document-Scale Upstream Pipeline (full tables)}
\label{app:doc-pipeline}
\setcounter{table}{0}\setcounter{figure}{0}

\Cref{fig:teaser} (main text) shows the three-stage schematic;
\Cref{tab:doc-pipeline-stages} below reports each stage's held-out
test performance and \Cref{tab:doc-pipeline-corpus} the per-category
instance counts. \Cref{fig:pipeline} gives the line-extractor cell
schematic referenced in \S\ref{sec:method}.

\begin{figure}[h]
\centering
\includegraphics[width=\linewidth]{figures/fig_pipeline.png}
\caption{Line-extractor cell schematic referenced from
\S\ref{sec:method}. Five cells run left to right: input chart $\to$
YOLO11 axis detector $\to$ YOLO11 legend masker $\to$ Mask2Former
foreground segmenter ($q{=}12$, $d{=}3$, $h{=}256$, MiT-B3 backbone,
multi-scale + flip TTA, threshold $\theta_{\mathrm{fg}}{=}0.25$) $\to$
$5$D LAB-$xy$ $k$-means separator with v4 post-processing
(morph-close $5$, dilate $3$, min-area $200$, $\Delta E$-gated
dilation). The gate $g(F)$ on the foreground fraction $\rho(F)$
reroutes the prediction to a scatter-only $k$-means branch when the
mask is nearly empty ($\rho(F){<}\tau_{\mathrm{fg}}{=}0.003$).}
\label{fig:pipeline}
\end{figure}

\begin{algorithm}[h]
\caption{The MatDeplot Stage~3 line-extractor on a single chart
image.}\label{alg:pipeline}
\begin{algorithmic}[1]
\Require image $I$; thresholds
$\theta_{\mathrm{fg}}{=}0.25$, $\tau_{\mathrm{fg}}{=}0.003$;
TTA scales $\mathcal{S}{=}\{512,768,1024\}$; spatial weight
$w_{xy}{=}0.5$.
\State $B_{\mathrm{axis}}, B_{\mathrm{legend}} \gets \textsc{YoloAxis}(I), \textsc{YoloLegend}(I)$;
fit calibration $\mathcal{C}$; build inside-plot mask $\mathcal{V}$.
\State $F \gets \mathrm{maxprob}\bigl(\{\textsc{M2F}_s(I), \textsc{M2F}_s(\mathrm{flip}\,I)\}_{s\in\mathcal{S}}\bigr)$.
\State $P \gets \{F\geq\theta_{\mathrm{fg}}\}\cap\mathcal{V}$;
$\rho \gets |P|/(H\cdot W)$.
\State $k \gets \textsc{KSource}(I)$
\Comment{cached Gemini one-shot, or silhouette in $\{2,\ldots,6\}$.}
\If{$\rho < \tau_{\mathrm{fg}}$}
  \State $\hat S \gets \textsc{KMeans}_k\bigl(\phi[\mathcal{V}\setminus\textsc{Otsu}(I)]\bigr)$;
  \texttt{used\_fallback}$\gets$True.
\Else
  \State $\hat S \gets \textsc{KMeans}_k\bigl(\phi[P]\bigr)$;
  apply v4 post-processing
  (morph-close~$5$, dilate~$3$, min-area~$200$, $\Delta E$-gated dilation);
  \texttt{used\_fallback}$\gets$False.
\EndIf
\State $\{\pi_i\} \gets \textsc{SkeletoniseAndOrder}(\hat S)$;
calibrate via $\mathcal{C}$.
\State \Return $\{\pi_i\}$ with per-curve confidence and
\texttt{used\_fallback} flag.
\end{algorithmic}
\end{algorithm}

\begin{table}[h]
\centering
\caption{Document-scale pipeline stages with held-out test
performance. Subfigure-detector results are mean$\pm$range of $5$
independent training runs; axis and legend detectors are
mean$\pm$range across the same $5$ runs on the line-extraction
test sets. Classifier accuracies are stratified by the eight
characterization categories on $n{=}1000$ images each, hand
verified.}
\label{tab:doc-pipeline-stages}
\scriptsize
\setlength{\tabcolsep}{2pt}
\begin{tabular}{rllrrl}
\toprule
Stage & Sub-step & Model & $|\text{Test}|$ & Headline metric & Source \\
\midrule
\multirow{4}{*}{\textbf{1}} & 1.1 PDF parsing            & MinerU                                  & 55{,}763 docs & --- & external                       \\
                            & 1.2 Subfigure splitting    & \textbf{YOLO11x} (this work)            & --- & mAP$_{50}$=98.4\%, mAP$_{50:95}$=93.2\% & Tab.~\ref{tab:subfig-detector} \\
                            & ($\alpha$ comparison)      & Mask R-CNN R-50                         & --- & mAP$_{50}$=96.5\%, mAP$_{50:95}$=76.9\% & Tab.~\ref{tab:subfig-detector} \\
                            & ($\alpha$ comparison)      & Mask R-CNN R-101                        & --- & mAP$_{50}$=97.4\%, mAP$_{50:95}$=78.6\% & Tab.~\ref{tab:subfig-detector} \\
\midrule
\textbf{2} & Hierarchical classifier & multimodal LLM (\texttt{figure\_class}) & 1{,}989 & top-1 acc=98.0\% (1950/1989) & Tab.~\ref{tab:doc-pipeline-classify} \\
\midrule
\multirow{4}{*}{\textbf{3}} & 3.1 Axis bbox detection    & YOLO11 (this work)                      & 108 instances & mAP$_{50}$=99.4\%, mAP$_{50:95}$=99.3\% & --- \\
                            & 3.2 Legend / clutter bbox  & YOLO11 (this work)                      & 594 instances & mAP$_{50}$=88.8\%, mAP$_{50:95}$=75.4\% & --- \\
                            & 3.3 Foreground segmentation & Mask2Former (this work)                & --- & --- & \S\ref{sec:method}                            \\
                            & 3.4 Per-curve separation   & 5D LAB-$xy$ $k$-means (\S\ref{sec:method}) & 204 manual & $\mathrm{LR}_{0.5}$=45.8\% & Tab.~\ref{tab:iou-headline} \\
\bottomrule
\end{tabular}
\end{table}

\begin{table}[h]
\centering
\caption{Subfigure detector head-to-head on the same $5$-run
train/val split. YOLO11x beats both Mask R-CNN backbones on every
metric and every individual run; the $+15$\,pp mAP$_{50:95}$ gap
is the dominant lever for the downstream classifier.}
\label{tab:subfig-detector}
\small
\begin{tabular}{lcccc}
\toprule
Detector & precision & recall & mAP$_{50}$ & mAP$_{50:95}$ \\
\midrule
\textbf{YOLO11x} (ours)     & \textbf{0.967}$\pm$0.005 & \textbf{0.962}$\pm$0.003 & \textbf{0.984}$\pm$0.002 & \textbf{0.932}$\pm$0.002 \\
Mask R-CNN R-50            & 0.838$\pm$0.011          & 0.822$\pm$0.013          & 0.965$\pm$0.011          & 0.769$\pm$0.020          \\
Mask R-CNN R-101           & 0.850$\pm$0.014          & 0.840$\pm$0.020          & 0.974$\pm$0.004          & 0.786$\pm$0.021          \\
\bottomrule
\end{tabular}
\end{table}

\begin{table}[h]
\centering
\caption{Multimodal-LLM classifier accuracy per characterization
category on $n{=}1000$ hand-verified subfigures per category.
\textbf{Failed} counts mis-routed subfigures; accuracy is
$1{-}\text{failed}/1000$.}
\label{tab:doc-pipeline-classify}
\small
\begin{tabular}{lccccc|cccc}
\toprule
& \multicolumn{5}{c|}{Structural characterisation} & \multicolumn{4}{c}{Electrochemical / kinetic} \\
& XRD & XPS & Raman & EIS & mean & GCD & CV & C-Rate & Rates \\
\midrule
failed (out of 1000) & 45 & 29 & 49 & 66 & 47.3 & 82 & 94 & 32 & 59 \\
accuracy             & 95.5\% & 97.1\% & 95.1\% & 93.4\% & 95.3\% & 91.8\% & 90.6\% & 96.8\% & 94.1\% \\
\bottomrule
\end{tabular}
\end{table}

\begin{table}[h]
\centering
\caption{Per-category line-chart instance counts produced by the
end-to-end document pipeline on the $55{,}763$-paper corpus.}
\label{tab:doc-pipeline-corpus}
\small
\begin{tabular}{lr|lr}
\toprule
Structural characterisation & count & Electrochemical / kinetic & count \\
\midrule
XRD               & 64{,}736  & GCD                  & 93{,}420 \\
XPS               & 63{,}296  & CV                   & 67{,}207 \\
Raman             & 18{,}629  & C-Rate               & 60{,}903 \\
EIS               & 36{,}383  & Rate-capability      & 35{,}581 \\
\midrule
sub-total         & 183{,}044 & sub-total            & 257{,}111 \\
\midrule
\multicolumn{3}{l}{\textbf{All categories}} &  \textbf{440{,}155} \\
\bottomrule
\end{tabular}
\end{table}

% ================================================================
\section{Mask2Former Architecture, Training-Data, and Loss Ablations}
\label{app:arch}
\setcounter{table}{0}\setcounter{figure}{0}

The default foreground segmenter ($q{=}12$, $d{=}3$, $h{=}256$,
MiT-B3 backbone) was selected from a 24-cell sweep over architecture
$(q,d,h)$, training-data slices ($20$k vs.\ $30$k vs.\ $40$k
synthetic charts, with or without $64{\times}$ real-image
oversampling), and three loss recipes:
\textbf{Recipe~A} (binary-IoU, our chosen training protocol),
\textbf{Recipe~B} (dice + focal, real-augment ablation), and
\textbf{Recipe~C} (multi-task instance + trajectory + style +
endpoint, plus BCE-only controls). All variants share an identical
optimiser configuration: AdamW, $100$ epochs, batch size $8$, image
size $640$, $\mathrm{lr}{=}1{\times}10^{-4}$ on the head and
$1{\times}10^{-5}$ on the MiT-B3 backbone, with linear warm-up over
$5$ epochs and cosine decay; trained on a single H100 80\,GB GPU.

\Cref{tab:arch-sweep} reports the best-epoch real validation
foreground IoU (the metric the checkpoint scheduler optimises) and
per-pixel F1 across all 24 trained variants. \Cref{tab:fg-ablation}
additionally re-uses four checkpoints under the full end-to-end
pipeline on the leakage-controlled leakage-clean $n{=}90$ hold --- the ranking
on validation IoU transfers monotonically to the downstream
DTW-$\Delta$ and CCR ranking. Three points stand out.

\begin{itemize}
\item \emph{Architecture knobs are sub-dominant.} Across the
$\{q{\in}\{8,10,12\},\,d{\in}\{3,4\},\,h{=}256\}$ sub-grid under
Recipe~A, the spread of best-epoch real-IoU is $\pm 0.005$ around
the chosen point; no $(q,d)$ gain exceeds the validation noise
floor.
\item \emph{Training-data scale dominates.} Moving the synthetic
slice from $20$k to $40$k images, or moving from synthetic-only to
$64{\times}$ real-image oversampling (synth-only vs.\
real$\times 64$ rows of \Cref{tab:arch-sweep}), shifts real-IoU by
$+0.05$--$+0.07$ --- an order of magnitude larger than the
architectural moves.
\item \emph{Multi-task loss design is mostly neutral.} Recipe~C
(multi-task) sits within $\pm 0.03$ of the simpler Recipe~A
binary-IoU; a BCE-only control is $-0.03$ behind the IoU loss;
removing the trajectory term (\emph{no-traj}) or shrinking the
dataset to the Recipe-B size cap (\emph{B-fair}, \emph{B-size})
loses $0.01$--$0.05$ IoU.
\end{itemize}

The combination ``Recipe~A architecture + $40$k synthetic + $182$
real charts, binary-IoU loss'' (the released
\texttt{q12\_d3\_h256} checkpoint) is therefore preferred not
because it dominates every alternative, but because it sits inside
the flat plateau on validation IoU while having the simplest loss
recipe and the cleanest end-to-end downstream profile
(\Cref{tab:fg-ablation}).

\begin{table}[h]
\centering
\caption{Architecture / data / loss sweep on the held-out real
validation split (24 trained checkpoints; best epoch per row).
\textsc{real\_iou} is the best-epoch real-image foreground IoU
(the metric the checkpoint scheduler optimises). Synthetic slice:
$20$k / $30$k / $40$k denote synthetic-data scale; ``raw'' / ``mask''
distinguishes two label-generation strategies on the $40$k slice.
\textsc{real$\times K$} = real images repeated $K{\times}$ during
training. Recipes: A = binary-IoU (default), B = dice $+$ focal,
C = multi-task (instance + trajectory + style + endpoint).}
\label{tab:arch-sweep}
\footnotesize
\setlength{\tabcolsep}{3pt}
\begin{tabular}{llcccp{2.5cm}p{2.4cm}r}
\toprule
Recipe & Variant & $q$ & $d$ & $h$ & slice & loss / notes & real\_iou \\
\midrule
\multirow{8}{*}{A}
 & q8\_d3\_h256\,(raw)                 &  8 & 3 & 256 & 40k raw  & binary-IoU                & 0.5498 \\
 & q10\_d3\_h256\,(raw)                & 10 & 3 & 256 & 40k raw  & binary-IoU                & 0.5531 \\
 & q10\_d4\_h256\,(raw)                & 10 & 4 & 256 & 40k raw  & binary-IoU                & 0.5531 \\
 & \textbf{q12\_d3\_h256\,(default)}   & 12 & 3 & 256 & 40k raw  & \textbf{binary-IoU}       & \textbf{0.5557} \\
 & q8\_d3\_h256\,(mask)                &  8 & 3 & 256 & 40k mask & binary-IoU                & 0.5552 \\
 & q10\_d3\_h256\,(mask)               & 10 & 3 & 256 & 40k mask & binary-IoU                & 0.5371 \\
 & q10\_d4\_h256\,(mask)               & 10 & 4 & 256 & 40k mask & binary-IoU                & 0.5481 \\
 & q12\_d3\_h256\,(mask)               & 12 & 3 & 256 & 40k mask & binary-IoU                & 0.5456 \\
\midrule
\multirow{9}{*}{B}
 & q8\_d3\_h256 (synth-only)           &  8 & 3 & 256 & 40k, synth-only          & dice+focal & 0.4864 \\
 & q8\_d3\_h256 (real$\times 64$)      &  8 & 3 & 256 & 40k, real$\times 64$     & dice+focal & 0.5279 \\
 & q10\_d3\_h256 (synth-only)          & 10 & 3 & 256 & 40k, synth-only          & dice+focal & 0.4837 \\
 & q10\_d3\_h256 (real$\times 64$)     & 10 & 3 & 256 & 40k, real$\times 64$     & dice+focal & 0.5371 \\
 & q10\_d4\_h256 (synth-only)          & 10 & 4 & 256 & 40k, synth-only          & dice+focal & 0.4858 \\
 & q10\_d4\_h256 (real$\times 64$)     & 10 & 4 & 256 & 40k, real$\times 64$     & dice+focal & 0.5367 \\
 & q12\_d3\_h256 (synth-only)          & 12 & 3 & 256 & 40k, synth-only          & dice+focal & 0.4828 \\
 & q12\_d3\_h256 (semantic-only)       & 12 & 3 & 256 & 40k, synth-only          & dice (semantic only) & 0.4844 \\
 & q12\_d3\_h256 (real$\times 64$)     & 12 & 3 & 256 & 40k, real$\times 64$     & dice+focal & 0.5484 \\
\midrule
\multirow{8}{*}{C}
 & main (seed 0)                        & 12 & 3 & 256 & 40k, real$\times 64$     & multi-task & 0.5322 \\
 & main (seed 1)                        & 12 & 3 & 256 & 40k, real$\times 64$     & multi-task & 0.5352 \\
 & queries=10                           & 10 & 3 & 256 & 40k                       & multi-task & 0.5093 \\
 & queries=30                           & 30 & 3 & 256 & 40k                       & multi-task & 0.5117 \\
 & no-trajectory loss                   & 12 & 3 & 256 & 40k                       & multi-task, $-$ trajectory & 0.5061 \\
 & BCE-only control                     & 12 & 3 & 256 & 40k                       & BCE-only       & 0.5269 \\
 & B-fair (size-matched)                & 12 & 3 & 256 & 40k (size-matched to B)  & multi-task & 0.5062 \\
 & B-size (size cap)                    & 12 & 3 & 256 & 40k (B-size cap)         & multi-task & 0.5243 \\
\bottomrule
\end{tabular}
\end{table}

\begin{table}[h]
\centering
\caption{End-to-end FG-checkpoint swap on the leakage-controlled
$n{=}90$ $n{=}90$ leakage-clean subset, holding YOLO axis, YOLO legend,
LAB-$k$-means separator, and skeleton-polyline extractor fixed
(no morph post-processing, no flip TTA, threshold $0.35$, scales
$\{640,1024\}$ --- a controlled FG-only ablation; not the
shape-mode or pixel-mode recipe of \Cref{tab:t1-all,tab:iou-headline}).
The DTW and CCR ranking is monotone with the real-IoU ranking of
\Cref{tab:arch-sweep}.}
\label{tab:fg-ablation}
\small
\begin{tabular}{lrrrrr}
\toprule
Checkpoint & DTW$\downarrow$ & CCR$\uparrow$ & Inst-F1$\uparrow$ &
             PQ$\uparrow$ & fb \\
\midrule
q12\_d3\_h256, $40$k synth $+$ $182$ real (default) & \textbf{19.92} & \textbf{0.898} & 0.650 & \textbf{0.204} & 0/90 \\
q10\_d3\_h256, $20$k synth $+$ $182$ real             & 23.78 & 0.876 & 0.672 & 0.191 & 1/90 \\
q15\_d3\_h256, $20$k synth $+$ $182$ real             & 23.27 & 0.877 & 0.690 & 0.178 & 1/90 \\
q10\_d6\_h256, $30$k synth $+$ $182$ real             & 22.80 & 0.880 & \textbf{0.691} & 0.202 & 0/90 \\
\bottomrule
\end{tabular}
\end{table}

% ================================================================
\section{Per-Stratum Tables and Auxiliary Sweeps}
\label{app:strata}\label{app:ablations}\label{app:agent-ready}
\label{app:ablations}
\setcounter{table}{0}\setcounter{figure}{0}

\Cref{tab:t1-all-full} extends \Cref{tab:t1-all} with the full
oracle-FG ceiling (M2--M6) and the UNet++/ResNeSt M4 row.
\Cref{tab:two-modes-app} reports the shape-mode and pixel-mode
counterparts of \Cref{tab:t1-all} side-by-side on the legacy $36$ set and
the $n{=}90$ hold (the two modes share the entire backbone and differ
only in the Cell~4 post-processing block; \S\ref{sec:method},
\emph{Two operating points}). \Cref{tab:t1-strata} reports the
per-stratum split for our \textsc{MatDeplot} (shape mode) and
\textsc{MatDeplot-NoFG} rows on the legacy $36$ set;
\Cref{tab:label-clean-90} is the matching per-stratum breakdown on the
leakage-free $n{=}90$ $n{=}90$ leakage-clean subset.
\Cref{tab:vlm-label-clean-90} is the live-API VLM head-to-head on
that same subset. \Cref{tab:thresh-sweep} is the fallback threshold
sweep, \Cref{tab:robust-sweep} the perturbation table.

\begin{table}[h]
\centering
\caption{Shape-mode vs.\ pixel-mode of MatDeplot's Cell~4 (the two
operating points share Cells~1--3 and Cell~5; only the
morph-close + min-area + LAB-$\Delta E$ dilation block differs).
On the legacy $36$ set the shape mode wins shape-regime metrics; the pixel
mode wins pixel-regime metrics on the $n{=}90$ hold. Sources:
\texttt{ours\_kmeans\_v1tta.csv}, \texttt{ours\_kmeans\_v4tta.csv},
\texttt{phaseC\_scales512\_full204\_metrics.csv} (shape on n=90,
no morph),
\texttt{phaseC\_flip\_full204\_metrics.csv} (pixel on n=90).}
\label{tab:two-modes-app}
\small
\setlength{\tabcolsep}{4pt}
\begin{tabular}{l|cccc|cc}
\toprule
& DTW$\downarrow$ & CCR$\uparrow$ & Inst-F1$\uparrow$ & Curve-PQ$\uparrow$
& $\mathrm{LR}_{0.5}\uparrow$ & $\mathrm{LR}_{0.8}\uparrow$ \\
\midrule
\multicolumn{7}{l}{\emph{the legacy $36$-image set}} \\
shape mode  & \textbf{14.7} & \textbf{0.964} & 0.520 & \textbf{0.178} & --- & --- \\
pixel mode  & 23.2          & 0.951          & \textbf{0.572} & 0.031 & --- & --- \\
\midrule
\multicolumn{7}{l}{\emph{leakage-clean $n{=}90$ hold ($n{=}90$, cached Gemini-K)}} \\
shape mode (gemK)  & 21.11          & 0.878          & 0.929          & \textbf{0.229} & --- & --- \\
pixel mode (gemK)  & \textbf{19.93} & \textbf{0.917} & \textbf{0.942} & 0.126 & \textbf{45.8\%} & \textbf{0.5\%} \\
\bottomrule
\end{tabular}\\
\footnotesize{LR columns are not reported for the legacy $36$ set because
the unified evaluator on that split records DTW/CCR/Inst-F1/Curve-PQ
only (the LR sweep is computed against the dedicated 5-px rasteriser
of \texttt{iou\_metric.py}, which lives on the $n{=}204$ side; we
leave that cell blank rather than back-fill across rasterisers).
The shape-mode LR cells on $n{=}90$ are also blank: shape mode's
mask is thinner than the $5$-px GT rasterisation, so an
\texttt{iou\_metric.py}-style strict-IoU sweep would be
floor-bound; the matching reference is the upper bound
$\mathrm{LR}_{\theta}\le \mathrm{PQ}/\theta$ ($0.171/0.5{=}34.2\%$).
The pixel mode is the LR-tuned recipe.}
\end{table}

\begin{table}[h]
\centering
\caption{Full oracle-FG ceiling (M2--M6) plus the M4 UNet++/ResNeSt
row, completing \Cref{tab:t1-all}. M$_i$ is the $i$-th separator
cell of the MatDeplot oracle-FG harness.}
\label{tab:t1-all-full}
\small
\begin{tabular}{lcccl}
\toprule
Method & DTW-$\Delta\downarrow$ & CCR@0.1$\uparrow$ &
Inst-F1$\uparrow$ & Family \\
\midrule
M2 (oracle FG) & 113.8 & 0.865 & 0.437 & oracle ceiling \\
M3 (oracle FG) & 102.3 & 0.963 & 0.434 & oracle ceiling \\
M4 (oracle FG) &  69.0 & 0.982 & 0.138 & oracle ceiling \\
M5 (oracle FG) & 127.1 & 0.782 & 0.437 & oracle ceiling \\
M6 (oracle FG) & 101.1 & 0.942 & 0.265 & oracle ceiling \\
M4 (UNet++/ResNeSt) & 93.2 & 0.916 & 0.279 & prior local-FG \\
\bottomrule
\end{tabular}
\end{table}

\begin{table}[h]
\centering
\caption{Per-stratum performance of our two pipeline variants
(computed from \texttt{experiments/results/ours\_kmeans\_m2f\{,\_fallback\}.csv}).
Per-image means; the $n{=}1$ L3 row is shown for completeness only.}
\label{tab:t1-strata}
\small
\begin{tabular}{lcccc|cccc}
\toprule
& \multicolumn{4}{c|}{\textsc{MatDeplot}} &
  \multicolumn{4}{c}{\textsc{MatDeplot-NoFG}} \\
Stratum              & DTW & CCR & F1 & PQ &
                       DTW & CCR & F1 & PQ \\
\midrule
ALL ($n{=}36$)                   & \textbf{15.9} & \textbf{0.932} & 0.601 & 0.233 &
                                   25.3 & 0.852 & 0.471 & 0.076 \\
L1 single ($n{=}17$)             & 8.2   & 0.988 & 0.471 & 0.237 &
                                   18.8  & 0.920 & 0.312 & 0.054 \\
L2 multi-disjoint ($n{=}6$)      & 22.4  & 0.937 & 0.658 & 0.319 &
                                   34.5  & 0.775 & 0.571 & 0.093 \\
L3 multi-overlap ($n{=}1$)       & 2.1   & 1.000 & 0.800 & 0.619 &
                                   14.5  & 0.853 & 0.667 & 0.274 \\
L4 same-colour ($n{=}6$)         & 23.9  & 0.885 & 0.744 & 0.166 &
                                   28.8  & 0.834 & 0.688 & 0.142 \\
L5 dense scatter ($n{=}6$)       & 25.4  & 0.805 & 0.571 & 0.140 &
                                   32.7  & 0.770 & 0.523 & 0.051 \\
\bottomrule
\end{tabular}
\end{table}

\begin{table}[h]
\centering
\caption{\textsc{MatDeplot} (pixel mode, v4$+$TTA: scales
$\{512,768,1024\}$, flip TTA, $\theta_{\mathrm{fg}}{=}0.25$, morph
post-processing, cached Gemini-K) on the leakage-free $n{=}90$
leakage-clean $n{=}90$ hold. Per-stratum means; \emph{fb} is the fraction
of images triggering the fallback. Source:
\texttt{ours\_v4tta\_gemK\_label.csv} filtered by
\texttt{label\_clean\_holdout.txt}.}
\label{tab:label-clean-90}
\small
\begin{tabular}{lrrrrrr}
\toprule
Stratum & $n$ & DTW$\downarrow$ & CCR$\uparrow$ & Inst-F1$\uparrow$ &
PQ$\uparrow$ & fb \\
\midrule
ALL                 & 90 & \textbf{19.93} & 0.917 & \textbf{0.942} & 0.126 & 0.00 \\
L1 single           & 29 &  9.89 & \textbf{0.992} & 0.937 & \textbf{0.195} & 0.00 \\
L2 two-line         & 28 & 28.94 & 0.858 & 0.962 & 0.060 & 0.00 \\
L3 few-line (3--4)  & 28 & 20.18 & 0.889 & 0.923 & 0.112 & 0.00 \\
L4 many-line ($\geq 5$)& 5 & 26.27 & 0.964 & \textbf{0.971} & 0.170 & 0.00 \\
\bottomrule
\end{tabular}
\end{table}

\begin{table}[h]
\centering
\caption{VLM head-to-head on the leakage-clean $n{=}90$ (leakage-free)
held-out complement ($n{=}90$), 2026-Q2 live API calls (GPT-5.4 /
Gemini-3.1-pro-preview / Qwen-3.5-plus via viviai.cc and
DashScope). Same GT pixel polylines as \Cref{tab:label-clean-90}.
Latency is wall-clock seconds per image including network
round-trip. $^\dagger$ Gemini-3.1-pro-preview emits zero trails on
every L4 ($\geq 5$ lines) image; reported DTW averages over the
$85$ images where Hungarian matching produced any pairs.}
\label{tab:vlm-label-clean-90}
\small
\begin{tabular}{lrrrrr}
\toprule
Method & DTW$\downarrow$ & CCR$\uparrow$ & Inst-F1$\uparrow$ &
         PQ$\uparrow$ & latency (s) \\
\midrule
\textsc{MatDeplot} (pixel, gemK) & \textbf{19.93} & 0.917 & \textbf{0.942} & \textbf{0.126} & $\approx 2.0$ \\
\midrule
GPT-5.4                          & 61.59 & 0.587 & 0.915 & $<10^{-3}$ & $\sim 20$ \\
Gemini-3.1-pro-preview$^\dagger$ & 73.42 & 0.210 & 0.343 & 0.006 & $\sim 30$ \\
Qwen-3.5-plus                    & 76.40 & 0.521 & 0.837 & $<10^{-3}$ & $\sim 107$ \\
Claude-Opus-4-6                  & \multicolumn{5}{c}{\emph{omitted: viviai.cc gateway has no Claude channel; cf.~\Cref{tab:t1-all}}} \\
\bottomrule
\end{tabular}
\end{table}

\begin{table}[h]
\centering
\caption{A2 fallback-trigger threshold sweep on the legacy $36$ set
($n{=}36$).}
\label{tab:thresh-sweep}
\small
\begin{tabular}{lcccc}
\toprule
$\tau_{\mathrm{fg}}$ & \# trigger fires & DTW-$\Delta$ & CCR & Inst.F1 \\
\midrule
0.001           & 0    & 15.88 & 0.932 & 0.601 \\
0.003 (default) & 0    & 15.88 & 0.932 & 0.601 \\
0.010           & 0    & 15.88 & 0.932 & 0.601 \\
0.030           & 19   & 20.23 & 0.896 & 0.521 \\
0.100           & 36   & 25.29 & 0.852 & 0.471 \\
0.300           & 36   & 25.29 & 0.852 & 0.471 \\
0.990 (forced)  & 36   & 25.29 & 0.852 & 0.471 \\
\bottomrule
\end{tabular}
\end{table}

\begin{table}[h]
\centering
\caption{A6 robustness: end-to-end metrics under three perturbations
applied to all 36 test images and re-scored by the same evaluator.}
\label{tab:robust-sweep}
\small
\begin{tabular}{lcccc}
\toprule
Perturbation & DTW-$\Delta$~$\downarrow$ & CCR@0.1~$\uparrow$ & Inst.F1~$\uparrow$ & Curve-PQ~$\uparrow$ \\
\midrule
Clean (reference)           & 15.88  & 0.932 & 0.601 & 0.233 \\
Rotation $+5^\circ$         & 35.33  & 0.831 & 0.537 & 0.000 \\
JPEG $q{=}30$               & 17.11  & 0.924 & 0.603 & 0.219 \\
Gaussian $\sigma{=}5$       & 14.09  & 0.920 & 0.568 & 0.254 \\
\bottomrule
\end{tabular}
\end{table}

\begin{table}[h]
\centering
\caption{Per-$K$ y-lookup accuracy under condition A (image only)
vs.\ B (image + ours extraction), Gemini-3.1-pro on $n{=}100$
images, $800$ questions, $\pm 10\%$ relative-error tolerance
(referenced by \S\ref{sec:downstream}).}
\label{tab:vlm-ylookup}
\small
\begin{tabular}{cccccc}
\toprule
$K$ & n images & n questions & acc (A) & acc (B) & $\Delta$ \\
\midrule
1   & 33  & 264 & 0.319 & \textbf{0.636} & $+0.317$ \\
2   & 28  & 224 & 0.188 & \textbf{0.308} & $+0.121$ \\
3   & 23  & 184 & 0.163 & \textbf{0.277} & $+0.114$ \\
4   & 8   & 64  & 0.156 & \textbf{0.266} & $+0.109$ \\
5   & 5   & 40  & 0.105 & 0.075          & $-0.030$ \\
$\geq 6$ & 3 & 24 & 0.083 & 0.083        & $\pm 0.000$ \\
\midrule
ALL & 100 & 800 & 0.215 & \textbf{0.388} & $\mathbf{+0.172}$ \\
\bottomrule
\end{tabular}
\end{table}

\section{Qualitative Comparison Across Extractors}
\label{app:qualitative}
\setcounter{table}{0}\setcounter{figure}{0}

\Cref{fig:reconstruction-strip} shows our pipeline's per-line cluster
masks alongside the GT polylines on four showcase images (one per
stratum), and \Cref{fig:appendix-comparison} extends the same four
images with hosted-VLM and DePlot columns for direct visual
comparison.

\begin{figure}[t]
\centering
\includegraphics[width=\linewidth]{figures/fig_reconstruction_strip.png}
\caption{Reconstruction strip on the four showcase images
(rows: $L_1$ single XPS, $L_2$ two-curve XRD, $L_3$ three-curve Raman,
$L_4$ six-curve XRD stack); columns: original chart, GT polylines
(green), our pipeline's per-line cluster masks (one colour per matched
curve). Each image was picked as the highest-pct$_{\mathrm{loc}}$
instance in its stratum subject to
$K_{\mathrm{pred}}{=}K_{\mathrm{gt}}$, on the MatDeplot configuration (\S\ref{sec:results}). For the same images under hosted VLMs
and DePlot, see \Cref{fig:appendix-comparison}.}
\label{fig:reconstruction-strip}
\end{figure}

\Cref{fig:appendix-comparison} renders the same four showcase
images one per stratum (the same picks as
\Cref{fig:reconstruction-strip}), but extends the per-row strip
to seven columns: original chart, GT polylines, our pipeline
(per-line cluster masks), and the four published-baseline
extractors (GPT-5.4, Gemini-3.1-pro, Qwen-3.5-plus, DePlot).
For the four hosted VLMs (GPT-5.4~\citep{openai_release_2026},
Gemini-3.1-pro~\citep{google_release_2026},
Qwen-3.5-plus~\citep{qwen_release_2026},
Claude Opus 4.6~\citep{anthropic_release_2026}) we re-ran the
extraction prompt (\texttt{experiments/iter/vlm\_extract\_showcase.py},
identical prompt to \Cref{tab:vlm-label-clean-90}) on these four
images and re-anchored the returned data-coordinate polylines into
pixel space using the chart's known axis ticks. For
DePlot~\citep{deplot} we ran \texttt{google/deplot}
(Pix2Struct~\citep{pix2struct} fine-tune for chart-to-table)
locally via \texttt{transformers}~\citep{huggingface}
(\texttt{experiments/iter/deplot\_extract\_showcase.py}); its
output is the structured table parsed back into per-series
polylines using \texttt{scipy}~\citep{scipy} numerical parsing.

Two robust qualitative findings hold across all four strata:
\emph{(a) Ours} produces a one-to-one match between cluster
colours and GT curves --- the cluster colours are placed
\emph{on} the GT curves, not adjacent to them, even on the L4
many-line stack. \emph{(b) The hosted VLMs, when they emit any
polyline at all (GPT-5.4 and Qwen-3.5-plus most consistently),
draw curves whose data-coordinates are roughly correct but whose
re-anchored pixel positions sit several axis-tick widths away
from the GT polyline}. This is the visual mechanism behind the
$\mathrm{LR}_\tau{\approx}0$ row in \Cref{tab:iou-headline}:
the VLM understands the chart's content (number of curves, peak
locations in data units) but cannot ground that understanding in
the source image's pixel coordinate system. \emph{(c) DePlot}
returns a single trail per chart regardless of $K_{\mathrm{gt}}$
in our scientific-chart corpus (\texttt{n\_predicted\_trails}=1
on every showcase), confirming the existing literature's note
that chart-to-table models are tuned on bar/pie charts and
truncate after a small number of points.

Claude Opus 4.6 is not shown: the viviai.cc gateway available to
us has no Claude channel that accepts tool-style structured
outputs for our prompt (cf.\ \Cref{tab:t1-all}); on the four
showcase images the API call returns an error
(\texttt{success=False}, \texttt{n\_predicted\_trails=0}), so
its column would be empty.

\begin{figure}[t]
\centering
\includegraphics[width=\linewidth]{figures/fig_appendix_comparison.png}
\caption{Qualitative comparison of pipeline extractions on the four showcase
images (rows: $L_1$ single XPS, $L_2$ two-curve XRD, $L_3$ three-curve Raman,
$L_4$ six-curve XRD stack). Columns: original / GT polylines (green) / Ours
(per-line cluster colours) / GPT-5.4 / Gemini-3.1-pro / Qwen-3.5-plus / DePlot.
For each VLM column the sub-caption reports
$K_{\mathrm{pred}}/K_{\mathrm{gt}}$. Hosted VLMs emit data-coordinate
polylines that we re-anchor into pixel space using the chart's known axis
ticks; the residual horizontal/vertical drift visible in those columns is the
visual signature of $\mathrm{LR}_\tau{\approx}0$ in \Cref{tab:iou-headline}.
DePlot collapses every chart to a single series. Source data:
\texttt{experiments/results/vlm\_showcase/<backend>/<image\_id>.json}.}
\label{fig:appendix-comparison}
\end{figure}

% ================================================================
\section{Per-Chart-Type Faithfulness Breakdown}
\label{app:per-chart}
\label{app:per-chart-type}
\setcounter{table}{0}\setcounter{figure}{0}

\Cref{tab:per-chart-type-lr} stratifies $\mathrm{LR}_\tau$ by the
seven materials-science chart types we observe in the
MatCurvs-204 corpus. Chart-type labels come from the
provenance-aware multimodal-LLM classifier of Stage~2
(Subfigure Classification, \Cref{fig:teaser}); $175/204$
images are typed, the remaining $30$ are flagged ``unknown''.
The single most important takeaway is that the
\emph{distribution of $\mathrm{LR}_\tau$ across chart types is
heavily uneven for our pipeline} (IR and Raman > XRD and XPS >
EIS, EXAFS, and DOS), but \emph{flat at zero for every hosted
VLM regardless of chart type}. Two implications: \emph{(i)} our
pipeline's failure modes are concentrated on rare and
visually-different chart types (EIS Nyquist plots: complex-plane
arcs, no peaks; EXAFS: wave-like oscillations; DOS: bipolar
electronic-structure curves) where the foreground model has $\le 5$
training images; \emph{(ii)} hosted VLMs do \emph{not} get any
chart type ``for free'', even XRD and Raman, where their training
data presumably contains the most exposure.

\begin{table}[h]
\centering
\caption{$\mathrm{LR}_\tau$ stratified by chart type on the full
$n{=}204$ MatCurvs-204 set; ``unknown'' = images not yet typed by
the upstream classifier (treated conservatively, not dropped). VLM
rows use the Curve-PQ-derived upper bound of \Cref{tab:iou-headline}.
Source: \texttt{paper\_tables/tab\_per\_chart\_type\_lr.csv}.}
\label{tab:per-chart-type-lr}
\small
\setlength{\tabcolsep}{4pt}
\begin{tabular}{llrrrrrrr}
\toprule
Method & Chart type & $n_{\mathrm{img}}$ & $n_{\mathrm{lines}}$ &
$\mathrm{LR}_{0.5}$ & $\mathrm{LR}_{0.6}$ &
$\mathrm{LR}_{0.7}$ & $\mathrm{LR}_{0.8}$ & mean IoU \\
\midrule
\multirow{8}{*}{\textsc{MatDeplot}}
& IR        & 9   & 15  & \textbf{86.7\%} & \textbf{73.3\%} & \textbf{46.7\%} & 0.0\%  & 0.637 \\
& Raman     & 36  & 74  & 71.6\%          & 48.7\%          & 16.2\%          & 1.4\%  & 0.558 \\
& XPS       & 36  & 67  & 65.7\%          & 52.2\%          & 25.4\%          & 1.5\%  & 0.555 \\
& XRD       & 81  & 180 & 64.4\%          & 44.4\%          & 16.7\%          & 1.7\%  & 0.537 \\
& EXAFS     & 5   & 14  & 7.1\%           & 0.0\%           & 0.0\%           & 0.0\%  & 0.247 \\
& DOS       & 3   & 6   & 0.0\%           & 0.0\%           & 0.0\%           & 0.0\%  & 0.305 \\
& EIS       & 4   & 12  & 0.0\%           & 0.0\%           & 0.0\%           & 0.0\%  & 0.231 \\
& unknown   & 30  & 87  & 25.3\%          & 16.1\%          & 1.1\%           & 0.0\%  & 0.324 \\
\midrule
GPT-5.4         & all types & 203 & 454 & $\leq 0.0$\% & $\leq 0.0$\% & $\leq 0.0$\% & $\leq 0.0$\% & 0.002 \\
Gemini-3.1-pro  & all types & 204 & 455 & $\leq 0.5$\% & $\leq 0.0$\% & $\leq 0.0$\% & $\leq 0.0$\% & 0.005 \\
Qwen-3.5-plus   & all types & 190 & 421 & $\leq 0.0$\% & $\leq 0.0$\% & $\leq 0.0$\% & $\leq 0.0$\% & 0.000 \\
\bottomrule
\end{tabular}
\end{table}

% ================================================================
\section{Paired Bootstrap: TTA Recipe Ablation}
\label{app:bootstrap-ci}
\setcounter{table}{0}\setcounter{figure}{0}

\Cref{tab:bootstrap-ci} reports paired-bootstrap CIs on the core
headline metrics for two nested comparisons on the leakage-clean $n{=}90$
hold ($10{,}000$ resamples, image-clustered).
\textbf{(a)~Flip-TTA isolated} (\texttt{phaseC\_flip} vs.\
\texttt{phaseC\_scales512}) shares scale set $(512,768,1024)$ and
threshold $0.25$ — only the flip-TTA toggle differs.
\textbf{(b)~iter12 vs.\ iter11} (\texttt{phaseC\_flip} vs.\
\texttt{phaseB\_q12\_d3}) is the cumulative recipe gain
(threshold, scale set, and flip TTA all applied at once).
The flip-only comparison shows recall ($\mathrm{cov}_{\mathrm{gt}}$)
$+2.98$\,pp $[+2.71, +3.27]$ and a small precision regression
$-1.45$\,pp $[-1.68, -1.22]$, both 95\% CIs excluding zero; the
cumulative comparison (the previous version of this table) attributed
$+9.07$\,pp recall to the same mechanism, three-fold over-counting.
We report both so readers can see how the recipe knobs compose.

\emph{Multiple-comparison note.} Each panel reports four CIs over
the same paired sample; with no FDR correction the family-wise
$\alpha$ is $\approx 0.185$. We mark CIs that exclude zero with
$\star$ (95\% comparison-wise) but readers should treat
``simultaneously-significant'' claims with appropriate caution.

\begin{table}[h]
\centering
\caption{Paired bootstrap on $\Delta$, leakage-clean $n{=}90$ hold,
$10{,}000$ image-clustered resamples per metric. Panel (a)
isolates flip-TTA contribution; panel (b) is cumulative iter12 vs
iter11. Comparison-wise 95\% CIs; family-wise $\alpha \approx 0.185$
without FDR. Source: \texttt{paper\_tables/tab\_bootstrap\_ci.csv}.}
\label{tab:bootstrap-ci}
\small
\begin{tabular}{lrrrrr}
\toprule
Metric & phaseC\_flip & baseline & $\Delta$ &
$95\%$ CI & sig. \\
\midrule
\multicolumn{6}{l}{\emph{(a)~Flip-TTA isolated: phaseC\_flip vs.\ phaseC\_scales512.}} \\
IoU$_g$                       & 0.6492 & 0.6415 & $+0.0077$ & $[+0.0047, +0.0105]$ & $\star$ \\
IoU$_m$                       & 0.4810 & 0.4729 & $+0.0081$ & $[+0.0026, +0.0164]$ & $\star$ \\
$\mathrm{cov}_{\mathrm{gt}}$  & 0.8306 & 0.8008 & $+0.0298$ & $[+0.0271, +0.0327]$ & $\star$ \\
$\mathrm{cov}_{\mathrm{pred}}$ & 0.7558 & 0.7703 & $-0.0145$ & $[-0.0168, -0.0122]$ & $\star$ (worse) \\
\midrule
\multicolumn{6}{l}{\emph{(b)~Cumulative iter12 vs.\ iter11: phaseC\_flip vs.\ phaseB\_q12\_d3.}} \\
IoU$_g$                       & 0.6492 & 0.6408 & $+0.0084$ & $[-0.0010, +0.0178]$ & ns      \\
IoU$_m$                       & 0.4810 & 0.4768 & $+0.0042$ & $[-0.0046, +0.0137]$ & ns      \\
$\mathrm{cov}_{\mathrm{gt}}$  & 0.8306 & 0.7398 & $+0.0907$ & $[+0.0859, +0.0957]$ & $\star$ \\
$\mathrm{cov}_{\mathrm{pred}}$ & 0.7558 & 0.8336 & $-0.0778$ & $[-0.0868, -0.0694]$ & $\star$ (worse) \\
\bottomrule
\end{tabular}\\
\footnotesize{$\star$ denotes a comparison-wise 95\% CI excluding 0.
The cumulative panel (b) overstates the contribution attributable to
flip TTA alone by combining three recipe changes (threshold $0.35\to
0.25$, scale set $\{640,1024\}\to\{512,768,1024\}$, flip on/off);
panel (a) is the clean flip-only ablation.}
\end{table}

% ================================================================
\section{Per-Backend Error Modes}
\label{app:errors}
\setcounter{table}{0}\setcounter{figure}{0}

\Cref{tab:errors} reports per-backend rates of three
distinct failure modes on the $n{=}90$ hold: (a) API/runtime
failure (\texttt{success}=False), (b) successful API call but
empty prediction (\texttt{success}=True $\land$
\texttt{n\_predicted\_trails}=0), and the median per-call
latency. The breakdown clarifies why our $\mathrm{LR}_\tau$
table reports VLMs at $\approx 0$ even when their
Inst-F1 is high: GPT-5.4 returns one curve per image but
miscalibrates the data-to-pixel map; Gemini-3.1-pro
\emph{succeeds} on every call but emits an empty trail list
$64.4\%$ of the time; Qwen-3.5-plus is $90\%$ available but
spends $128$\,s per call; Claude Opus 4.6 fails on
$\approx 100\%$ of calls (gateway content blocking; we report
the partial $n{=}2$ row for transparency).

\begin{table}[h]
\centering
\caption{Error-mode breakdown on the $n{=}90$ hold. Median
latency is the $p_{50}$ across successful calls. Source:
\texttt{paper\_tables/tab\_errors.csv}.}
\label{tab:errors}
\small
\begin{tabular}{lrrrr}
\toprule
Method & $n$ & success & empty-pred & median $p_{50}$ (s) \\
\midrule
GPT-5.4                          & 90 & 1.000 & 0.000 & 19.2 \\
Gemini-3.1-pro                   & 90 & 1.000 & \textbf{0.644} & 33.1 \\
Qwen-3.5-plus                    & 90 & 0.900 & 0.000 & 128.1 \\
Claude Opus 4.6                  & \phantom{0}2 & 0.000 & ---   & 1.2  \\
\textbf{MatDeplot} & 90 & \textbf{1.000} & 0.000 & \textbf{1.5}  \\
\bottomrule
\end{tabular}
\end{table}

\begin{table}[h]
\centering
\caption{Agent-readiness scorecard (full table; App.~\ref{app:agent-ready}).
MatDeplot uses one cached Gemini K-call per unique image (\$0.012
amortised) and so loses the local-execution point; the FG-off
ablation drops the call entirely and reaches $7/7$. $p_{50}$
measured under the unified evaluator on the legacy $36$ set.}
\label{tab:agent-ready}
\small
\setlength{\tabcolsep}{4pt}
\begin{tabular}{lccccccccc}
\toprule
Method & schema & conf. & abstain & pixel & local &
$p_{50}(\mathrm{s})$ & cost$(\$)$ & Score \\
\midrule
\textsc{MatDeplot}      & 1.00 & \checkmark & \checkmark & \checkmark & --- & 1.5 & 0.012 & 6/7 \\
\textsc{MatDeplot-NoFG} & 1.00 & \checkmark & \checkmark & \checkmark & \checkmark & 1.5 & 0 & \textbf{7/7} \\
GPT-5.4                 & 0.92 & --- & --- & --- & --- & 44 & 0.009 & 2/7 \\
Claude Opus 4.6         & 1.00 & --- & --- & --- & --- & 16 & 0.118 & 1/7 \\
Gemini 3.1 Pro          & 0.83 & --- & --- & --- & --- & 35 & 0.016 & 0/7 \\
Qwen 3.5 Plus           & 0.83 & --- & --- & --- & --- & 126 & 0.012 & 0/7 \\
\bottomrule
\end{tabular}
\end{table}

\begin{table}[h]
\centering
\caption{Loop-scale agent-readiness measurements on MatCurvs-204
(every $204$ test images run end-to-end through each system; latency
in seconds, single A100 batch~$1$; cost in USD/image; abstain
fraction is per emitted curve). MatDeplot emits $1{,}035$ curves
across the $204$ images and abstains on $2$ ($0.19\%$) at the
default conservative $\tau_{\mathrm{abs}}{=}0.30$; per-curve
confidences are released as part of the JSON record so downstream
agents can re-threshold.}
\label{tab:agent-ready-detail}
\small
\setlength{\tabcolsep}{4pt}
\begin{tabular}{lrrrrrr}
\toprule
Method & schema-valid & empty-pred & $p_{50}$ & $p_{95}$ & $p_{99}$ & cost \\
\midrule
\textbf{MatDeplot}     & \textbf{1.000} & 0.000          & \textbf{1.5} & \textbf{2.6} & \textbf{3.7} & \textbf{0.012} \\
GPT-5.4                & 0.922          & 0.000          & 19.2         & 47.1         & 89.4         & 0.009 \\
Gemini-3.1-pro         & 1.000          & \textbf{0.644} & 33.1         & 64.0         & 91.2         & 0.016 \\
Qwen-3.5-plus          & 0.833          & 0.000          & 128.1        & 198.0        & 277.5        & 0.012 \\
Claude Opus 4.6        & 1.000          & ---            & 16.0         & ---          & ---          & 0.118 \\
\bottomrule
\end{tabular}
\end{table}

% ================================================================
\section{Raman Peak Detection: Detail by Peak Count}
\label{app:peak-error-detail}
\setcounter{table}{0}\setcounter{figure}{0}

\Cref{tab:raman-peak} reports the headline numbers for the Raman
pilot of \S\ref{sec:downstream} (axis-mapping ablation: v1/v2 use
the YOLO axis box for $x$ bounds, v3 uses the cluster mask's
left/right extrema), and \Cref{tab:peak-error-detail} stratifies
the v3 ($n{=}60$) results by GT peak count. The headline medians
($\mathrm{recall}{=}1.00$, $\mathrm{MAE}{=}2.58$\,cm$^{-1}$,
Pearson${=}0.955$) hold uniformly across peak counts, including
spectra with $9$+ peaks where peak detection is hardest. Median
precision rises slightly with peak count, consistent with the
intensity-threshold post-filter being more conservative when there
are more peaks to choose from.

\begin{table}[h]
\centering
\caption{Raman peak detection from extracted curves on $n{=}60$
spectra (medians; precision is reported only at v3 because it
became meaningful only after the axis-mapping correction).
Instrument resolution ${\sim}4$\,cm$^{-1}$.}
\label{tab:raman-peak}
\small
\begin{tabular}{lcccc}
\toprule
Variant & recall ($\uparrow$) & precision ($\uparrow$) &
MAE cm$^{-1}$ ($\downarrow$) & Pearson ($\uparrow$) \\
\midrule
v1 (axis\_box, n=30)        & 0.40  & 0.43 & 9.64 & 0.35 \\
v2 (axis\_box, n=60)        & 0.50  & ---  & 9.98 & 0.32 \\
\textbf{v3 (cluster x, n=60)} & \textbf{1.00} & \textbf{0.71}
                          & \textbf{2.58} & \textbf{0.96} \\
\bottomrule
\end{tabular}
\end{table}

\begin{table}[h]
\centering
\caption{Raman peak detection from extracted curves
(v3 = cluster-x bounds; n=60), stratified by GT peak count per
spectrum. Source:
\texttt{paper\_tables/tab\_peak\_error\_detail.csv}.}
\label{tab:peak-error-detail}
\small
\begin{tabular}{lrrrrr}
\toprule
Peak count bin & $n$ & median recall & median precision &
median MAE (cm$^{-1}$) & median Pearson \\
\midrule
1--4 peaks   & 12 & 1.000 & 0.667 & 2.53 & 0.903 \\
5--8 peaks   & 14 & 1.000 & 0.690 & 3.41 & 0.950 \\
9--12 peaks  & 34 & 1.000 & 0.769 & 2.55 & 0.958 \\
\midrule
ALL          & 60 & \textbf{1.000} & 0.714 & \textbf{2.58} & \textbf{0.955} \\
\bottomrule
\end{tabular}
\end{table}

% ================================================================
\section{Reproducibility Pointers}
\label{app:repro}
\setcounter{table}{0}\setcounter{figure}{0}

\begin{itemize}
\item \textbf{Raw per-image metrics}:
\texttt{experiments/results/ours\_kmeans\_m2f.csv} (36 rows,
MatDeplot pipeline);
\texttt{ours\_kmeans\_m2f\_fallback.csv} (36 rows, FG-off ablation).
\item \textbf{Cross-method benchmark snapshot}:
\texttt{results\_snapshot/per\_method\_metrics.csv} ($526$ rows
over $16$ methods $\times$ up to $36$ legacy $36$-set images).
\item \textbf{Per-image stratum labels}:
\texttt{results\_snapshot/image\_strata.json}.
\item \textbf{Pre-trained Mask2Former checkpoint}:
\texttt{models/mask2former\_q12\_d3\_h256\_best.pt} (released under
CC-BY-SA-4.0).
\item \textbf{Unified evaluator}:
\texttt{matdeplot.eval.unified\_metrics} (released with this paper
under MIT).
\item \textbf{Commit hash and seed}: documented in the released
repository's \texttt{REPRO.md}.
\end{itemize}

% ================================================================
\section{Leakage Audit (MatCurvs-204 vs.\ FG Training Set)}
\label{app:leakage}
\setcounter{table}{0}\setcounter{figure}{0}

The leakage audit summarised here is the in-PDF self-contained
version of \texttt{LABEL\_LEAKAGE\_FIX.md} (released alongside the
code), responding to the contamination concern raised in
\S\ref{sec:limits}\,(b).

\paragraph{Audit procedure.}
The Mask2Former-q12 FG model is trained on a corpus of $182$
hand-annotated real charts plus $40$k synthetic charts (\S\ref{sec:method}).
The MatCurvs-204 evaluation set is hand-annotated on a different
roll-up of the same upstream corpus, so file-level overlap with the
training set is possible. We run a SHA-256 audit at the bytes-of-image
level: for every MatCurvs-204 \texttt{.jpg}, we fingerprint the
pixel buffer and check membership against the training root's
fingerprint set. $114/204$ ($55.9\%$) MatCurvs-204 images have
identical SHA-256 to a training image; the remaining $90$ images
are byte-distinct from every training image.

\paragraph{The leakage-clean $n{=}90$ hold.}
We construct \texttt{label\_clean\_holdout.txt} from the $90$
byte-disjoint complement and restrict every headline number
(\S\ref{sec:results}, \Cref{tab:iou-headline}) to that subset. The
empirical leakage premium of the pixel-mode recipe is
$\mathrm{LR}_{0.5}{=}54.7\%$ (contaminated $n{=}204$) vs.\
$45.8\%$ (clean $n{=}90$) --- a $9$\,pp drop --- and
$\mathrm{LR}_{0.8}{=}1.1\%$ vs.\ $0.5\%$ ($-0.6$\,pp). Tighter
thresholds compress the gap because tight-IoU is harder to
``memorise'' through near-duplicates. We report the
clean-$n{=}90$ numbers as headline throughout.

\paragraph{What the audit cannot rule out.}
SHA-256 is exact-match; an image cropped or compressed differently
would evade the audit. We additionally checked perceptual-hash
nearness (pHash distance $\leq 4$ on a $64$-bit hash) against the
training root and found $0$ further hits beyond the SHA-256 set,
consistent with the test annotation pipeline being a deduplicated
subset of a different file dump. Two confounds remain: (i)~the
authors of papers in the test set \emph{may} have published other
figures used in training (but the FG model trains on chart images
not on paper-author bundles, so this transfer is at most
distribution-level); (ii)~the cached Gemini-3.1-pro $k$-source
(\S\ref{sec:method}, ``$k$-source: hybrid call'') was queried on
every MatCurvs-204 image including the $n{=}90$ hold, an
information channel disclosed in \S\ref{sec:limits}\,(c).

\paragraph{Reproduction.}
The audit script \texttt{experiments/iter/audit\_leakage.py} reads
the training-root manifest from \texttt{data/training\_manifest.json}
and the $n{=}204$ image list from \texttt{label/}, computes the SHA
fingerprints, and writes \texttt{label\_clean\_holdout.txt} (the
$90$-image complement) and \texttt{label\_overlap.txt} (the $114$
contaminated). The audit is deterministic; rerun is one command.

% ================================================================
\section{VLM-as-Judge Benchmark}
\label{app:vlm-judge}
\setcounter{table}{0}\setcounter{figure}{0}

We use four hosted VLMs as judges of the MatDeplot pixel-mode
extraction (v4$+$TTA $+$ cached Gemini-K, the recipe of
\Cref{tab:iou-headline}) on every MatCurvs-204 image. The
pipeline-extracted polylines are rendered as a high-contrast
overlay on the original chart with each curve labelled
\texttt{\#0}\,..\,\texttt{\#K{-}1}, and the prompt asks for a
per-curve PASS/FAIL plus an image-level pass-iff-all aggregate
(\texttt{vlm\_reviewer.py}; prompt verbatim in
\texttt{REVIEW\_PROMPT}). Three judges produced valid verdicts
(Qwen-VL-Max, GPT-5.4, Gemini-3.1-pro-preview); Claude Opus~4.6
returned $503$ \texttt{model\_not\_found} on every call through
the viviai gateway and is omitted. Total API spend $\sim$\$6;
full verdicts in \texttt{judge\_ours\_v4tta\_gemK\_label.csv}.

\paragraph{Per-judge pass rate.}
\Cref{tab:judge-passrate} reports the image-level pass rate on the
full $n{=}204$ set and on the leakage-clean $n{=}90$ hold.

\begin{table}[h]
\centering
\caption{Per-judge image-level pass rate. Source:
\texttt{judge\_passrate\_ours\_v4tta\_label.csv}.}
\label{tab:judge-passrate}
\small
\begin{tabular}{lrrrrr}
\toprule
Judge & $n_{204}$ & pass$_{204}$ & $n_{90}$ & pass$_{90}$ & median latency (s) \\
\midrule
Qwen-VL-Max                 & 204 & 0.544 &  90 & 0.422 &   4.4 \\
GPT-5.4                     & 204 & 0.093 &  90 & 0.111 &   5.1 \\
Gemini-3.1-pro-preview      & 204 & 0.108 &  90 & 0.100 &  16.1 \\
Claude Opus 4.6             & \multicolumn{5}{c}{\emph{omitted: viviai gateway returned $503$ on every call (no Claude channel)}} \\
\bottomrule
\end{tabular}
\end{table}

\paragraph{Cross-judge agreement.}
\Cref{tab:judge-kappa} reports Cohen's $\kappa$ between every pair
of judges over the common image set (judges that errored on a
specific image are excluded from that pair only). Values $\geq 0.6$
indicate substantial agreement; $\geq 0.8$ near-perfect.

\begin{table}[h]
\centering
\caption{Cohen's $\kappa$ on image-level PASS verdicts. Lower
triangle redundant; diagonal is $1.0$ by definition. Source:
\texttt{judge\_kappa\_ours\_v4tta\_label.csv}.}
\label{tab:judge-kappa}
\small
\begin{tabular}{lcccc}
\toprule
& Qwen-VL-Max & GPT-5.4 & Gemini-3.1-pro & Claude Opus 4.6 \\
\midrule
Qwen-VL-Max & 1.000    & 0.085    & 0.111    & ---      \\
GPT-5.4  &          & 1.000    & 0.052    & ---      \\
Gemini-3.1-pro-preview &          &          & 1.000    & ---      \\
Claude Opus 4.6 &          &          &          & 1.000    \\
\bottomrule
\end{tabular}
\end{table}

\paragraph{Validity: judge flags vs.\ GT-derived LR.}
For each judge we compute the mean GT-derived $\mathrm{LR}_{0.5}$
on images the judge \emph{passed} vs.\ \emph{failed}; a higher
$\mathrm{LR}_{0.5}$ on the pass set validates that the judge's
flag tracks real reconstruction quality, while
$\Delta{=}\mathrm{LR}_{0.5}^{\mathrm{pass}}-
\mathrm{LR}_{0.5}^{\mathrm{fail}}$ measures the strength of the
correlation.

\begin{table}[h]
\centering
\caption{Judge image-level verdicts vs.\ GT-derived $\mathrm{LR}_{0.5}$
on the leakage-clean $n{=}90$ hold. Positive $\Delta$ means
judge-passed images really are better-reconstructed.}
\label{tab:judge-vs-lr}
\small
\begin{tabular}{lcccc}
\toprule
Judge & $n^{\mathrm{pass}}$ & $\mathrm{LR}_{0.5}^{\mathrm{pass}}$ &
        $\mathrm{LR}_{0.5}^{\mathrm{fail}}$ & $\Delta$ \\
\midrule
Qwen-VL-Max                 & 38 & 0.766 & 0.279 &  0.487 \\
GPT-5.4                     & 10 & 0.800 & 0.445 &  0.355 \\
Gemini-3.1-pro-preview      &  9 & 0.889 & 0.440 &  0.449 \\
Claude Opus 4.6             & --- & --- & --- & --- \\
\bottomrule
\end{tabular}
\end{table}

\paragraph{Consensus rule for Cell~6 abstain.}
Claude Opus~4.6 returned $503$ on every call (viviai gateway has no
Claude channel; \S\ref{sec:limits}), so the consensus operates over
the $3$ effective judges (Qwen-VL-Max, GPT-5.4, Gemini-3.1-pro).
\Cref{tab:judge-consensus} sweeps the FAIL-vote threshold; the
$\geq 3$-of-$3$ \emph{unanimous-fail} rule abstains on $47/90$
images ($52.2\%$) and yields the strongest correlation with
GT-derived $\mathrm{LR}_{0.5}$ ($0.770$ on the kept subset vs.\
$0.223$ on the abstained subset, $\Delta{=}{+}0.546$). The looser
$\geq 2$-of-$3$ rule abstains on $78/90$ ($86.7\%$, overflagging)
because individual judges disagree substantially (pairwise
$\kappa{=}0.05$--$0.11$, \Cref{tab:judge-kappa}). We therefore
recommend the unanimous-fail rule as the deployment default.
Per-curve abstain flags from individual judges are released
alongside the image-level verdicts in
\texttt{judge\_ours\_v4tta\_gemK\_label.csv}.

\begin{table}[h]
\centering
\caption{Cell-6 consensus-abstain sweep on the $n{=}90$ hold:
abstain rate vs.\ GT-derived $\mathrm{LR}_{0.5}$ on the kept and
abstained subsets, for thresholds of $\geq 1, 2, 3$ judges flagging
FAIL. Source: \texttt{judge\_consensus\_ours\_v4tta\_label.csv}.}
\label{tab:judge-consensus}
\small
\begin{tabular}{ccccc}
\toprule
$\geq$ FAIL votes & abstain rate & $\mathrm{LR}_{0.5}^{\mathrm{kept}}$ &
                   $\mathrm{LR}_{0.5}^{\mathrm{abst.}}$ & $\Delta$ \\
\midrule
1 & 0.978 & 0.500 & 0.484 &  0.016 \\
\textbf{2 (default)} & 0.867 & 0.917 & 0.418 &  0.499 \\
3 & 0.522 & 0.770 & 0.223 &  0.546 \\
\bottomrule
\end{tabular}
\end{table}

\newpage
\input{checklist.tex}

\end{document}
